% concatenated from iclr2026_conference.tex, math_commands.tex

\documentclass{article} % For LaTeX2e
\usepackage{iclr2026_conference,times}
\iclrfinalcopy

% Optional math commands from https://github.com/goodfeli/dlbook_notation.
%%%%% NEW MATH DEFINITIONS %%%%%

\usepackage{amsmath,amsfonts,bm}

% Mark sections of captions for referring to divisions of figures
\newcommand{\figleft}{{\em (Left)}}
\newcommand{\figcenter}{{\em (Center)}}
\newcommand{\figright}{{\em (Right)}}
\newcommand{\figtop}{{\em (Top)}}
\newcommand{\figbottom}{{\em (Bottom)}}
\newcommand{\captiona}{{\em (a)}}
\newcommand{\captionb}{{\em (b)}}
\newcommand{\captionc}{{\em (c)}}
\newcommand{\captiond}{{\em (d)}}

% Highlight a newly defined term
\newcommand{\newterm}[1]{{\bf #1}}


% Figure reference, lower-case.
\def\figref#1{figure~\ref{#1}}
% Figure reference, capital. For start of sentence
\def\Figref#1{Figure~\ref{#1}}
\def\twofigref#1#2{figures \ref{#1} and \ref{#2}}
\def\quadfigref#1#2#3#4{figures \ref{#1}, \ref{#2}, \ref{#3} and \ref{#4}}
% Section reference, lower-case.
\def\secref#1{section~\ref{#1}}
% Section reference, capital.
\def\Secref#1{Section~\ref{#1}}
% Reference to two sections.
\def\twosecrefs#1#2{sections \ref{#1} and \ref{#2}}
% Reference to three sections.
\def\secrefs#1#2#3{sections \ref{#1}, \ref{#2} and \ref{#3}}
% Reference to an equation, lower-case.
\def\eqref#1{equation~\ref{#1}}
% Reference to an equation, upper case
\def\Eqref#1{Equation~\ref{#1}}
% A raw reference to an equation---avoid using if possible
\def\plaineqref#1{\ref{#1}}
% Reference to a chapter, lower-case.
\def\chapref#1{chapter~\ref{#1}}
% Reference to an equation, upper case.
\def\Chapref#1{Chapter~\ref{#1}}
% Reference to a range of chapters
\def\rangechapref#1#2{chapters\ref{#1}--\ref{#2}}
% Reference to an algorithm, lower-case.
\def\algref#1{algorithm~\ref{#1}}
% Reference to an algorithm, upper case.
\def\Algref#1{Algorithm~\ref{#1}}
\def\twoalgref#1#2{algorithms \ref{#1} and \ref{#2}}
\def\Twoalgref#1#2{Algorithms \ref{#1} and \ref{#2}}
% Reference to a part, lower case
\def\partref#1{part~\ref{#1}}
% Reference to a part, upper case
\def\Partref#1{Part~\ref{#1}}
\def\twopartref#1#2{parts \ref{#1} and \ref{#2}}

\def\ceil#1{\lceil #1 \rceil}
\def\floor#1{\lfloor #1 \rfloor}
\def\1{\bm{1}}
\newcommand{\train}{\mathcal{D}}
\newcommand{\valid}{\mathcal{D_{\mathrm{valid}}}}
\newcommand{\test}{\mathcal{D_{\mathrm{test}}}}

\def\eps{{\epsilon}}


% Random variables
\def\reta{{\textnormal{$\eta$}}}
\def\ra{{\textnormal{a}}}
\def\rb{{\textnormal{b}}}
\def\rc{{\textnormal{c}}}
\def\rd{{\textnormal{d}}}
\def\re{{\textnormal{e}}}
\def\rf{{\textnormal{f}}}
\def\rg{{\textnormal{g}}}
\def\rh{{\textnormal{h}}}
\def\ri{{\textnormal{i}}}
\def\rj{{\textnormal{j}}}
\def\rk{{\textnormal{k}}}
\def\rl{{\textnormal{l}}}
% rm is already a command, just don't name any random variables m
\def\rn{{\textnormal{n}}}
\def\ro{{\textnormal{o}}}
\def\rp{{\textnormal{p}}}
\def\rq{{\textnormal{q}}}
\def\rr{{\textnormal{r}}}
\def\rs{{\textnormal{s}}}
\def\rt{{\textnormal{t}}}
\def\ru{{\textnormal{u}}}
\def\rv{{\textnormal{v}}}
\def\rw{{\textnormal{w}}}
\def\rx{{\textnormal{x}}}
\def\ry{{\textnormal{y}}}
\def\rz{{\textnormal{z}}}

% Random vectors
\def\rvepsilon{{\mathbf{\epsilon}}}
\def\rvtheta{{\mathbf{\theta}}}
\def\rva{{\mathbf{a}}}
\def\rvb{{\mathbf{b}}}
\def\rvc{{\mathbf{c}}}
\def\rvd{{\mathbf{d}}}
\def\rve{{\mathbf{e}}}
\def\rvf{{\mathbf{f}}}
\def\rvg{{\mathbf{g}}}
\def\rvh{{\mathbf{h}}}
\def\rvu{{\mathbf{i}}}
\def\rvj{{\mathbf{j}}}
\def\rvk{{\mathbf{k}}}
\def\rvl{{\mathbf{l}}}
\def\rvm{{\mathbf{m}}}
\def\rvn{{\mathbf{n}}}
\def\rvo{{\mathbf{o}}}
\def\rvp{{\mathbf{p}}}
\def\rvq{{\mathbf{q}}}
\def\rvr{{\mathbf{r}}}
\def\rvs{{\mathbf{s}}}
\def\rvt{{\mathbf{t}}}
\def\rvu{{\mathbf{u}}}
\def\rvv{{\mathbf{v}}}
\def\rvw{{\mathbf{w}}}
\def\rvx{{\mathbf{x}}}
\def\rvy{{\mathbf{y}}}
\def\rvz{{\mathbf{z}}}

% Elements of random vectors
\def\erva{{\textnormal{a}}}
\def\ervb{{\textnormal{b}}}
\def\ervc{{\textnormal{c}}}
\def\ervd{{\textnormal{d}}}
\def\erve{{\textnormal{e}}}
\def\ervf{{\textnormal{f}}}
\def\ervg{{\textnormal{g}}}
\def\ervh{{\textnormal{h}}}
\def\ervi{{\textnormal{i}}}
\def\ervj{{\textnormal{j}}}
\def\ervk{{\textnormal{k}}}
\def\ervl{{\textnormal{l}}}
\def\ervm{{\textnormal{m}}}
\def\ervn{{\textnormal{n}}}
\def\ervo{{\textnormal{o}}}
\def\ervp{{\textnormal{p}}}
\def\ervq{{\textnormal{q}}}
\def\ervr{{\textnormal{r}}}
\def\ervs{{\textnormal{s}}}
\def\ervt{{\textnormal{t}}}
\def\ervu{{\textnormal{u}}}
\def\ervv{{\textnormal{v}}}
\def\ervw{{\textnormal{w}}}
\def\ervx{{\textnormal{x}}}
\def\ervy{{\textnormal{y}}}
\def\ervz{{\textnormal{z}}}

% Random matrices
\def\rmA{{\mathbf{A}}}
\def\rmB{{\mathbf{B}}}
\def\rmC{{\mathbf{C}}}
\def\rmD{{\mathbf{D}}}
\def\rmE{{\mathbf{E}}}
\def\rmF{{\mathbf{F}}}
\def\rmG{{\mathbf{G}}}
\def\rmH{{\mathbf{H}}}
\def\rmI{{\mathbf{I}}}
\def\rmJ{{\mathbf{J}}}
\def\rmK{{\mathbf{K}}}
\def\rmL{{\mathbf{L}}}
\def\rmM{{\mathbf{M}}}
\def\rmN{{\mathbf{N}}}
\def\rmO{{\mathbf{O}}}
\def\rmP{{\mathbf{P}}}
\def\rmQ{{\mathbf{Q}}}
\def\rmR{{\mathbf{R}}}
\def\rmS{{\mathbf{S}}}
\def\rmT{{\mathbf{T}}}
\def\rmU{{\mathbf{U}}}
\def\rmV{{\mathbf{V}}}
\def\rmW{{\mathbf{W}}}
\def\rmX{{\mathbf{X}}}
\def\rmY{{\mathbf{Y}}}
\def\rmZ{{\mathbf{Z}}}

% Elements of random matrices
\def\ermA{{\textnormal{A}}}
\def\ermB{{\textnormal{B}}}
\def\ermC{{\textnormal{C}}}
\def\ermD{{\textnormal{D}}}
\def\ermE{{\textnormal{E}}}
\def\ermF{{\textnormal{F}}}
\def\ermG{{\textnormal{G}}}
\def\ermH{{\textnormal{H}}}
\def\ermI{{\textnormal{I}}}
\def\ermJ{{\textnormal{J}}}
\def\ermK{{\textnormal{K}}}
\def\ermL{{\textnormal{L}}}
\def\ermM{{\textnormal{M}}}
\def\ermN{{\textnormal{N}}}
\def\ermO{{\textnormal{O}}}
\def\ermP{{\textnormal{P}}}
\def\ermQ{{\textnormal{Q}}}
\def\ermR{{\textnormal{R}}}
\def\ermS{{\textnormal{S}}}
\def\ermT{{\textnormal{T}}}
\def\ermU{{\textnormal{U}}}
\def\ermV{{\textnormal{V}}}
\def\ermW{{\textnormal{W}}}
\def\ermX{{\textnormal{X}}}
\def\ermY{{\textnormal{Y}}}
\def\ermZ{{\textnormal{Z}}}

% Vectors
\def\vzero{{\bm{0}}}
\def\vone{{\bm{1}}}
\def\vmu{{\bm{\mu}}}
\def\vtheta{{\bm{\theta}}}
\def\va{{\bm{a}}}
\def\vb{{\bm{b}}}
\def\vc{{\bm{c}}}
\def\vd{{\bm{d}}}
\def\ve{{\bm{e}}}
\def\vf{{\bm{f}}}
\def\vg{{\bm{g}}}
\def\vh{{\bm{h}}}
\def\vi{{\bm{i}}}
\def\vj{{\bm{j}}}
\def\vk{{\bm{k}}}
\def\vl{{\bm{l}}}
\def\vm{{\bm{m}}}
\def\vn{{\bm{n}}}
\def\vo{{\bm{o}}}
\def\vp{{\bm{p}}}
\def\vq{{\bm{q}}}
\def\vr{{\bm{r}}}
\def\vs{{\bm{s}}}
\def\vt{{\bm{t}}}
\def\vu{{\bm{u}}}
\def\vv{{\bm{v}}}
\def\vw{{\bm{w}}}
\def\vx{{\bm{x}}}
\def\vy{{\bm{y}}}
\def\vz{{\bm{z}}}

% Elements of vectors
\def\evalpha{{\alpha}}
\def\evbeta{{\beta}}
\def\evepsilon{{\epsilon}}
\def\evlambda{{\lambda}}
\def\evomega{{\omega}}
\def\evmu{{\mu}}
\def\evpsi{{\psi}}
\def\evsigma{{\sigma}}
\def\evtheta{{\theta}}
\def\eva{{a}}
\def\evb{{b}}
\def\evc{{c}}
\def\evd{{d}}
\def\eve{{e}}
\def\evf{{f}}
\def\evg{{g}}
\def\evh{{h}}
\def\evi{{i}}
\def\evj{{j}}
\def\evk{{k}}
\def\evl{{l}}
\def\evm{{m}}
\def\evn{{n}}
\def\evo{{o}}
\def\evp{{p}}
\def\evq{{q}}
\def\evr{{r}}
\def\evs{{s}}
\def\evt{{t}}
\def\evu{{u}}
\def\evv{{v}}
\def\evw{{w}}
\def\evx{{x}}
\def\evy{{y}}
\def\evz{{z}}

% Matrix
\def\mA{{\bm{A}}}
\def\mB{{\bm{B}}}
\def\mC{{\bm{C}}}
\def\mD{{\bm{D}}}
\def\mE{{\bm{E}}}
\def\mF{{\bm{F}}}
\def\mG{{\bm{G}}}
\def\mH{{\bm{H}}}
\def\mI{{\bm{I}}}
\def\mJ{{\bm{J}}}
\def\mK{{\bm{K}}}
\def\mL{{\bm{L}}}
\def\mM{{\bm{M}}}
\def\mN{{\bm{N}}}
\def\mO{{\bm{O}}}
\def\mP{{\bm{P}}}
\def\mQ{{\bm{Q}}}
\def\mR{{\bm{R}}}
\def\mS{{\bm{S}}}
\def\mT{{\bm{T}}}
\def\mU{{\bm{U}}}
\def\mV{{\bm{V}}}
\def\mW{{\bm{W}}}
\def\mX{{\bm{X}}}
\def\mY{{\bm{Y}}}
\def\mZ{{\bm{Z}}}
\def\mBeta{{\bm{\beta}}}
\def\mPhi{{\bm{\Phi}}}
\def\mLambda{{\bm{\Lambda}}}
\def\mSigma{{\bm{\Sigma}}}

% Tensor
\DeclareMathAlphabet{\mathsfit}{\encodingdefault}{\sfdefault}{m}{sl}
\SetMathAlphabet{\mathsfit}{bold}{\encodingdefault}{\sfdefault}{bx}{n}
\newcommand{\tens}[1]{\bm{\mathsfit{#1}}}
\def\tA{{\tens{A}}}
\def\tB{{\tens{B}}}
\def\tC{{\tens{C}}}
\def\tD{{\tens{D}}}
\def\tE{{\tens{E}}}
\def\tF{{\tens{F}}}
\def\tG{{\tens{G}}}
\def\tH{{\tens{H}}}
\def\tI{{\tens{I}}}
\def\tJ{{\tens{J}}}
\def\tK{{\tens{K}}}
\def\tL{{\tens{L}}}
\def\tM{{\tens{M}}}
\def\tN{{\tens{N}}}
\def\tO{{\tens{O}}}
\def\tP{{\tens{P}}}
\def\tQ{{\tens{Q}}}
\def\tR{{\tens{R}}}
\def\tS{{\tens{S}}}
\def\tT{{\tens{T}}}
\def\tU{{\tens{U}}}
\def\tV{{\tens{V}}}
\def\tW{{\tens{W}}}
\def\tX{{\tens{X}}}
\def\tY{{\tens{Y}}}
\def\tZ{{\tens{Z}}}


% Graph
\def\gA{{\mathcal{A}}}
\def\gB{{\mathcal{B}}}
\def\gC{{\mathcal{C}}}
\def\gD{{\mathcal{D}}}
\def\gE{{\mathcal{E}}}
\def\gF{{\mathcal{F}}}
\def\gG{{\mathcal{G}}}
\def\gH{{\mathcal{H}}}
\def\gI{{\mathcal{I}}}
\def\gJ{{\mathcal{J}}}
\def\gK{{\mathcal{K}}}
\def\gL{{\mathcal{L}}}
\def\gM{{\mathcal{M}}}
\def\gN{{\mathcal{N}}}
\def\gO{{\mathcal{O}}}
\def\gP{{\mathcal{P}}}
\def\gQ{{\mathcal{Q}}}
\def\gR{{\mathcal{R}}}
\def\gS{{\mathcal{S}}}
\def\gT{{\mathcal{T}}}
\def\gU{{\mathcal{U}}}
\def\gV{{\mathcal{V}}}
\def\gW{{\mathcal{W}}}
\def\gX{{\mathcal{X}}}
\def\gY{{\mathcal{Y}}}
\def\gZ{{\mathcal{Z}}}

% Sets
\def\sA{{\mathbb{A}}}
\def\sB{{\mathbb{B}}}
\def\sC{{\mathbb{C}}}
\def\sD{{\mathbb{D}}}
% Don't use a set called E, because this would be the same as our symbol
% for expectation.
\def\sF{{\mathbb{F}}}
\def\sG{{\mathbb{G}}}
\def\sH{{\mathbb{H}}}
\def\sI{{\mathbb{I}}}
\def\sJ{{\mathbb{J}}}
\def\sK{{\mathbb{K}}}
\def\sL{{\mathbb{L}}}
\def\sM{{\mathbb{M}}}
\def\sN{{\mathbb{N}}}
\def\sO{{\mathbb{O}}}
\def\sP{{\mathbb{P}}}
\def\sQ{{\mathbb{Q}}}
\def\sR{{\mathbb{R}}}
\def\sS{{\mathbb{S}}}
\def\sT{{\mathbb{T}}}
\def\sU{{\mathbb{U}}}
\def\sV{{\mathbb{V}}}
\def\sW{{\mathbb{W}}}
\def\sX{{\mathbb{X}}}
\def\sY{{\mathbb{Y}}}
\def\sZ{{\mathbb{Z}}}

% Entries of a matrix
\def\emLambda{{\Lambda}}
\def\emA{{A}}
\def\emB{{B}}
\def\emC{{C}}
\def\emD{{D}}
\def\emE{{E}}
\def\emF{{F}}
\def\emG{{G}}
\def\emH{{H}}
\def\emI{{I}}
\def\emJ{{J}}
\def\emK{{K}}
\def\emL{{L}}
\def\emM{{M}}
\def\emN{{N}}
\def\emO{{O}}
\def\emP{{P}}
\def\emQ{{Q}}
\def\emR{{R}}
\def\emS{{S}}
\def\emT{{T}}
\def\emU{{U}}
\def\emV{{V}}
\def\emW{{W}}
\def\emX{{X}}
\def\emY{{Y}}
\def\emZ{{Z}}
\def\emSigma{{\Sigma}}

% entries of a tensor
% Same font as tensor, without \bm wrapper
\newcommand{\etens}[1]{\mathsfit{#1}}
\def\etLambda{{\etens{\Lambda}}}
\def\etA{{\etens{A}}}
\def\etB{{\etens{B}}}
\def\etC{{\etens{C}}}
\def\etD{{\etens{D}}}
\def\etE{{\etens{E}}}
\def\etF{{\etens{F}}}
\def\etG{{\etens{G}}}
\def\etH{{\etens{H}}}
\def\etI{{\etens{I}}}
\def\etJ{{\etens{J}}}
\def\etK{{\etens{K}}}
\def\etL{{\etens{L}}}
\def\etM{{\etens{M}}}
\def\etN{{\etens{N}}}
\def\etO{{\etens{O}}}
\def\etP{{\etens{P}}}
\def\etQ{{\etens{Q}}}
\def\etR{{\etens{R}}}
\def\etS{{\etens{S}}}
\def\etT{{\etens{T}}}
\def\etU{{\etens{U}}}
\def\etV{{\etens{V}}}
\def\etW{{\etens{W}}}
\def\etX{{\etens{X}}}
\def\etY{{\etens{Y}}}
\def\etZ{{\etens{Z}}}

% The true underlying data generating distribution
\newcommand{\pdata}{p_{\rm{data}}}
% The empirical distribution defined by the training set
\newcommand{\ptrain}{\hat{p}_{\rm{data}}}
\newcommand{\Ptrain}{\hat{P}_{\rm{data}}}
% The model distribution
\newcommand{\pmodel}{p_{\rm{model}}}
\newcommand{\Pmodel}{P_{\rm{model}}}
\newcommand{\ptildemodel}{\tilde{p}_{\rm{model}}}
% Stochastic autoencoder distributions
\newcommand{\pencode}{p_{\rm{encoder}}}
\newcommand{\pdecode}{p_{\rm{decoder}}}
\newcommand{\precons}{p_{\rm{reconstruct}}}

\newcommand{\laplace}{\mathrm{Laplace}} % Laplace distribution

\newcommand{\E}{\mathbb{E}}
\newcommand{\Ls}{\mathcal{L}}
\newcommand{\R}{\mathbb{R}}
\newcommand{\emp}{\tilde{p}}
\newcommand{\lr}{\alpha}
\newcommand{\reg}{\lambda}
\newcommand{\rect}{\mathrm{rectifier}}
\newcommand{\softmax}{\mathrm{softmax}}
\newcommand{\sigmoid}{\sigma}
\newcommand{\softplus}{\zeta}
\newcommand{\KL}{D_{\mathrm{KL}}}
\newcommand{\Var}{\mathrm{Var}}
\newcommand{\standarderror}{\mathrm{SE}}
\newcommand{\Cov}{\mathrm{Cov}}
% Wolfram Mathworld says $L^2$ is for function spaces and $\ell^2$ is for vectors
% But then they seem to use $L^2$ for vectors throughout the site, and so does
% wikipedia.
\newcommand{\normlzero}{L^0}
\newcommand{\normlone}{L^1}
\newcommand{\normltwo}{L^2}
\newcommand{\normlp}{L^p}
\newcommand{\normmax}{L^\infty}

\newcommand{\parents}{Pa} % See usage in notation.tex. Chosen to match Daphne's book.

\DeclareMathOperator*{\argmax}{arg\,max}
\DeclareMathOperator*{\argmin}{arg\,min}

\DeclareMathOperator{\sign}{sign}
\DeclareMathOperator{\Tr}{Tr}
\let\ab\allowbreak


\usepackage{hyperref}
\usepackage{url}
\usepackage[ruled,vlined,linesnumbered]{algorithm2e}
\SetKwInput{KwIn}{Input}
\SetKwInput{KwOut}{Output}
\usepackage{svg}
\usepackage{stmaryrd}
\usepackage{xcolor}
\usepackage{amsmath,amssymb,amsthm}
\newtheorem{definition}{Definition}
\newtheorem{theorem}{Theorem}[section]
\renewcommand{\thetheorem}{\arabic{theorem}}
\newtheorem{remark}{Remark}
\newtheorem{lemma}{Lemma}
\newtheorem{corollary}{Corollary}
\newtheorem{assumption}{Assumption}
\newtheorem{proposition}{Proposition}
\usepackage{enumitem}
\usepackage{multirow}

\usepackage{subcaption}
\usepackage{graphicx}
\usepackage{booktabs}



%\title{Federated Learning with Biased Compression via Step-Ahead Partial Error Feedback}
%\title{SA-PEF: A Unified Scheme for Communication-Efficient Federated Learning under Heterogeneous Data}
\title{SA-PEF: Step-Ahead Partial Error Feedback for Efficient Federated Learning}

% Authors must not appear in the submitted version. They should be hidden
% as long as the \iclrfinalcopy macro remains commented out below.
% Non-anonymous submissions will be rejected without review.

\author{Dawit Kiros Redie\textsuperscript{*},
Reza Arablouei\textsuperscript{$\dagger$},
and Stefan Werner\textsuperscript{*,$\ddagger$}\\
\textsuperscript{*}Department of Electronic Systems, Norwegian University of Science and Technology (NTNU), Norway\\
\textsuperscript{$\ddagger$}Department of Information and Communications Engineering, Aalto University, Finland\\
\textsuperscript{$\dagger$}CSIRO’s Data61, Pullenvale, QLD 4069, Australia\\
\texttt{\{dawit.k.redie, stefan.werner\}@ntnu.no, reza.arablouei@data61.csiro.au}
}

% The \author macro works with any number of authors. There are two commands
% used to separate the names and addresses of multiple authors: \And and \AND.
%
% Using \And between authors leaves it to \LaTeX{} to determine where to break
% the lines. Using \AND forces a linebreak at that point. So, if \LaTeX{}
% puts 3 of 4 authors names on the first line, and the last on the second
% line, try using \AND instead of \And before the third author name.

\newcommand{\fix}{\marginpar{FIX}}
\newcommand{\new}{\marginpar{NEW}}

%\iclrfinalcopy % Uncomment for camera-ready version, but NOT for submission.
\begin{document}


\maketitle

\begin{abstract}

Biased gradient compression with error feedback (EF) reduces communication in federated learning (FL), but under non-IID data, the residual error can decay slowly, causing gradient mismatch and stalled progress in the early rounds. We propose step-ahead partial error feedback (SA-PEF), which integrates step-ahead (SA) correction with partial error feedback (PEF). SA-PEF recovers EF when the step-ahead coefficient $\alpha=0$ and step-ahead EF (SAEF) when $\alpha=1$. For non-convex objectives and $\delta$-contractive compressors, we establish a second-moment bound and a residual recursion that guarantee convergence to stationarity under heterogeneous data and partial client participation. The resulting rates match standard non-convex Fed-SGD guarantees up to constant factors, achieving $O((\eta\,\eta_0TR)^{-1})$ convergence to a variance/heterogeneity floor with a fixed inner step size. Our analysis reveals a step-ahead-controlled residual contraction $\rho_r$ that explains the observed acceleration in the early training phase. To balance SAEF’s rapid warm-up with EF’s long-term stability, we select $\alpha$ near its theory-predicted optimum. Experiments across diverse architectures and datasets show that SA-PEF consistently reaches target accuracy faster than EF.

\end{abstract}


\section{Introduction}

Modern large‐scale machine learning increasingly relies on distributed computation, where both data and compute are spread across many devices. Federated learning (FL) enables model training in this setting without centralizing raw data, enhancing privacy and scalability under heterogeneous client distributions~\citep{mcmahan2017communication,kairouz2021advances}. In FL, a potentially vast population of clients collaborates to train a shared model $w\in\mathbb{R}^d$ by solving 
\begin{equation}\label{eq:global-obj}
w^\star \in \arg\min_{w\in\mathbb{R}^d} f(w) := \frac{1}{K}\sum_{k=1}^K f_k(w),
\qquad
f_k(w):=\mathbb{E}_{z\sim\mathcal{D}_k}\bigl[\ell(w;z)\bigr],
\end{equation}
where $\mathcal{D}_k$ is the (potentially heterogeneous) data distribution at client $k$, $\ell(\cdot)$ is a sample loss (often nonconvex), and $K$ is the number of clients.  
In each synchronous FL round, the server broadcasts the current global model to a subset of clients. These clients perform several steps of stochastic gradient descent (SGD) on their local data and return updates to the server, which aggregates them to form the next global iterate \citep{huang2022lower, wang2022unified,li2024low}. 

%FL aims to use rich, yet heterogeneous, data distributed across the network. Despite these benefits, it faces significant communication bottlenecks because model updates are high dimensional, especially for modern models with millions or billions of parameters \citep{reisizadeh2020fedpaq,kim2024spafl,islamov2025safeef}. This challenge has motivated strong interest in communication-efficient algorithms \citep{seide2014onebit,alistarh2017qsgd,stich2018sparsified,beznosikov2023biased,bao2025efskip}. Another central challenge is statistical heterogeneity, that is, when client data are typically non-IID. When multiple local training steps are taken, local models drift toward minimizing their own losses, which can hinder rapid convergence of the global model \citep{karimireddy2020scaffold,li2023analysis}.
Although FL leverages rich distributed data, it faces two key challenges. The first challenge is communication bottlenecks. Model updates are typically high-dimensional, with millions or even billions of parameters, which makes uplink bandwidth a major constraint \citep{reisizadeh2020fedpaq,kim2024spafl,islamov2025safeef}. This has spurred extensive work on communication-efficient algorithms, including quantization \citep{seide2014onebit,alistarh2017qsgd}, sparsification \citep{stich2018sparsified}, and biased compression with error feedback \citep{beznosikov2023biased,bao2025efskip}. The second challenge is statistical heterogeneity. When client data are non-IID, multiple local updates can cause client models to drift toward minimizing their own objectives, slowing or even destabilizing global convergence \citep{karimireddy2020scaffold,li2023analysis}.  

To reduce communication, many methods compress client-server messages using quantization \citep{alistarh2017qsgd}, sparsification \citep{lin2018deep}, or sketching \citep{rothchild2020fetchsgd}. Compressors may be unbiased (e.g., Rand-$k$ \citep{wangni2018gradient, stich2018sparsified}) or biased (e.g., signSGD \citep{bernstein2018signsgd}, Top-$k$ \citep{lin2018deep}), with the latter often delivering better accuracy-communication trade-offs at a given bit budget \citep{beznosikov2023biased}. However, naive use of biased compression can introduce a persistent bias, leading to slow or even divergent training \citep{beznosikov2023biased, li2023analysis}. Error feedback (EF) addresses this problem by storing past compression errors and injecting them into the next update before compression \citep{seide2014onebit}. This mechanism cancels the compressor bias and restores convergence guarantees comparable to those of uncompressed SGD, assuming standard smoothness and appropriately chosen stepsizes \citep{karimireddy2019error, bao2025efskip}.

EF in federated settings faces two important limitations. First, under highly non-IID data, the residual can align with client-specific gradient directions, inducing cross-client gradient mismatch and slowing early progress \citep{hsu2020saef, karimireddy2020scaffold}. Second, EF retains residual mass until fully transmitted. Once the residual norm is small, communication rounds may be wasted transmitting stale, low-magnitude coordinates rather than fresh gradient signal \citep{li2023analysis}.

Step-ahead EF (SAEF) \citep{xu2021step} mitigates the first issue by \emph{previewing} the residual: before local SGD, each client shifts its model by the current error and optimizes from that offset. This strategy often yields a strong warm-up, as the residual is injected in full every round. However, in FL regimes with non-IID data, multiple local steps, aggressive compression, and partial participation, the full step-ahead variant exhibits late-stage plateaus and larger gradient-mismatch spikes compared to EF.
%under noisy compression, SAEF can amplify variance in the model and hurt late-stage stability, sometimes plateauing above the baseline. 
Moreover, prior analysis~\citep{xu2021step} has largely focused on classical distributed optimization, leaving open whether one can \emph{systematically} combine SAEF's fast initial progress with the long-term stability of EF in federated settings with local steps and data heterogeneity.

\textbf{Our approach: SA-PEF.}
We propose \emph{step-ahead partial error feedback (SA-PEF)}, by introducing a tunable coefficient $\alpha_r\in[0,1]$ and shifting only a \emph{fraction} of the residual ($w_{r+\frac12}=w_r-\alpha_r e_r$) while carrying the remainder through standard EF. This partial shift provides several benefits:
\begin{itemize}[leftmargin=*,itemsep=0pt,topsep=0pt,parsep=0pt]
  \item \emph{Early acceleration with reduced noise.} A moderate (or decaying) $\alpha_r$ removes most early-round mismatch, while the injected noise automatically diminishes as $\|e_r\|$ shrinks.
  \item \emph{Tighter theoretical recursions.} We establish a residual contraction
        \(
        \rho_r=(1-\tfrac1\delta)\bigl[2(1-\alpha_r)^2+24\,\alpha_r^2(\eta_rLT)^2\bigr],
        \)
        which, for small $s_r=\eta_rLT$, is strictly smaller than the EF value $2(1-\tfrac1\delta)$. The local-drift constants also improve with $\alpha_r$.
  \item \emph{Graceful interpolation.} SA-PEF reduces to EF when $\alpha_r=0$ (maximal stability) and to SAEF when $\alpha_r=1$ (maximal jump-start), allowing practitioners to adapt to different data and communication regimes.
\end{itemize}

\textbf{Contributions.}
\begin{itemize}[leftmargin=*,itemsep=0pt,topsep=0pt,parsep=0pt]
\item \emph{Algorithm.} We introduce SA-PEF, a lightweight drop-in variant of Local-SGD with biased compression and a tunable step-ahead coefficient $\alpha_r$, compatible with any $\delta$-contractive compressor.
%\item \emph{Theory.} We derive new bounds specific to SA-PEF: (i) a local-model drift inequality, (ii) a second-moment bound for the shifted average update $\E\|\alpha_r\bar e_r-\bar g_r\|^2$, and (iii) a residual recursion with contraction factor $\rho_r$. These results recover EF constants at $\alpha_r=0$ and reveal how step-ahead modifies both drift and residual memory.
%\item \emph{Rates.} We prove a one-round descent inequality with explicit stepsize coupling, yielding nonconvex stationarity guarantees of order $O\big((\eta\,\eta_0TR)^{-1}\big)$ with constant inner stepsize. The dependence on compression appears only through $(1-1/\delta)$ and $\rho_{\max}$, matching prior compressed-FL analyses.
\item \emph{Theory.} We provide a convergence analysis of SA-PEF. 
Our results include new inequalities for drift, second moments, and residual recursion, which recover EF constants at $\alpha_r=0$ and quantify how step-ahead alters drift and residual memory. From these, we derive nonconvex stationarity guarantees of order $O\big((\eta\,\eta_0 T R)^{-1}\big)$, with compression dependence appearing only through $(1-1/\delta)$ and $\rho_{\max}$, in line with prior compressed-FL work.
\end{itemize}


\section{Related work}

\paragraph{Error feedback and compressed optimization.}
Error feedback (EF) was first introduced as a practical heuristic for 1-bit SGD~\citep{seide2014onebit} and later formalized as a \emph{memory} mechanism in sparsified or biased SGD \citep{stich2018sparsified,karimireddy2019error}. By accumulating the compression residual and injecting it into subsequent updates, EF restores descent directions and admits convergence guarantees for broad classes of \emph{contractive (possibly biased)} compressors. These results include linear rates in the strongly convex setting and standard stationary-point guarantees in the nonconvex case \citep{gorbunov2020linearly,beznosikov2023biased}. More recently, EF21~\citep{richtarik2021ef21} and its extensions~\citep{fatkhullin2025ef21bells} provide a modern error-feedback framework for compressing full gradients (or gradient differences) at a shared iterate in synchronized data-parallel training with $T{=}1$, achieving clean contraction guarantees and removing the error floor. However, these works assume no local steps and no client drift. 
% In the federated \emph{local-SGD} setting with $T>1$ local steps, partial participation, and biased contractive compressors, Fed-EF~\citep{li2023analysis} adapts classical EF and is the natural EF-style baseline that we extend with SA-PEF. 
A complementary line of work replaces residual memory with control variates (global gradient estimators), as in DIANA and MARINA \citep{mishchenko2019distributed,gorbunov2021marina}, which reduce or remove compressor bias without maintaining a full residual vector. \emph{EControl} \citep{gao2024econtrol} regulates the strength of the feedback signal and fuses residual and estimator updates into a single compressed message, providing fast convergence under arbitrary contractive compressors and heterogeneous data.

\paragraph{Local updates in FL.}
Local or periodic averaging (a.k.a. Local-SGD) reduces communication rounds by performing $T>1$ local steps between synchronizations \citep{stich2018local}. While effective in homogeneous settings, non-IID data induces \emph{client drift}, where model trajectories diverge across clients, degrading both convergence speed and final accuracy. Several approaches mitigate drift while retaining the communication savings of local updates. \emph{Proximal regularization} (FedProx) stabilizes local objectives by penalizing deviation from the current global model \citep{li2020federated}. \emph{Control variates} (SCAFFOLD) estimate and correct the client-specific gradient bias caused by heterogeneity, yielding tighter convergence with multiple local steps \citep{karimireddy2020scaffold}. \textit{Dynamic regularization} (FedDyn) further aligns local stationary points with the global objective via a round-wise correction term, improving robustness on highly non-IID data \citep{acar2021federated}.

\paragraph{Compression with local updates.}
Combining local updates with message compression compounds communication savings but also amplifies distortions from both local drift and compression error. FedPAQ \citep{reisizadeh2020fedpaq} performs $T>1$ local steps and transmits quantized model deltas at synchronization points, exposing explicit trade-offs among the local period, stepsizes, and quantization accuracy. QSparse-Local-SGD \citep{basu2019qsparse} extends this to contractive compressors, transmitting Top-$k$ updates after $T$ local steps. While achieving significant traffic reduction, it also reveals that aggressive sparsity can destabilize convergence. CSER \citep{xie2020cser} mitigates this with \emph{error reset}, which immediately injects the residual back into the local model to restore stability under high compression. 
In the federated local-SGD setting with partial participation and biased compression, Fed-EF~\citep{li2023analysis} provides a first nonconvex analysis of classical EF and serves as the EF-style backbone that SA-PEF builds upon. 
On the control-variate side, 
%SCAFFOLD introduces variance-reduced local updates without compression. 
Scaffnew/ProxSkip \citep{mishchenko2022proxskip} is a more recent local-training method in the SCAFFOLD family, using probabilistic local updates to achieve theoretical acceleration. However, it still relies on full-precision exchanges and no inherent compression. CompressedScaffnew \citep{condat2022provably} extends this mechanism with quantization, TAMUNA \citep{condat2023tamuna} further handles partial client participation, and LoCoDL \citep{condat2025locodl} generalizes the analysis to arbitrary unbiased compressors. These methods primarily establish accelerated convergence in (strongly) convex settings. In parallel, SCALLION/SCAFCOM \citep{huang2024stochastic} combines SCAFFOLD-style control variates with compression and, in SCAFCOM, local momentum to handle heterogeneity and partial participation, at the cost of additional per-client state.

\paragraph{Step-ahead error feedback.}
SAEF \citep{xu2021step} addresses \emph{gradient mismatch}, i.e., the discrepancy between the model used for gradient computation and the model actually updated when delayed residuals are applied. SAEF performs a \emph{preview} shift of the model using the residual before local SGD and augments this with occasional \emph{error averaging} across workers. Although this reduces mismatch and accelerates early progress, the analysis is developed for classical distributed settings with single-step synchronization and \emph{bounded-gradient} assumptions, and does not cover federated regimes with multiple local steps, non-IID data, or biased compressors. Moreover, error averaging requires extra coordination and communication, which is often impractical in cross-device FL.

%Two questions remain open for FL with local updates and biased compression: (i) how to control gradient mismatch \emph{without} extra communication (e.g., error averaging), and (ii) how to couple a step-ahead shift with EF so that early acceleration does not compromise late-stage stability under high compression and non-IID data. Our work closes this gap by (a) performing a controlled step-ahead shift, (b) retaining only a fraction of the residual each round (partial EF), and (c) developing a contraction-based analysis that links mismatch to a measurable residual energy and yields nonconvex stationary-point bounds under data heterogeneity.
Despite progress, it remains unclear how to (i) control gradient mismatch in federated settings with local steps and biased compression \emph{without} extra communication, or (ii) combine step-ahead correction with EF to balance early acceleration and long-term stability. Our work closes this gap by introducing SA-PEF, which performs a controlled step-ahead shift with partial residual retention on top of Fed-EF and provides a contraction-based analysis yielding nonconvex guarantees under heterogeneous data.

\section{Proposed algorithm}

{\small
\begin{algorithm}[ht]
\DontPrintSemicolon
\caption{\small Step‑Ahead Partial Error‑Feedback (SA‑PEF) for Efficient FL}
\label{alg:sa_pef}

\KwIn{total communication rounds $R$, client number $K$, stepsizes $\{\eta_r\}_{r=0}^{R-1}$, step‑ahead schedule $\{\alpha_r\}_{r=0}^{R-1}$ with $0\le\alpha_r\le 1$, compressor $\mathcal{C}(\cdot)$, initial model $w_0$, local step number $T$}

\ForEach{client $k\!=\!1,\dots,K$ {in parallel}}{
  $w_{0}^{(k)} \gets w_0$; \quad $e_{0}^{(k)} \gets 0$\;
}

\For{$r \gets 0$ \KwTo $R-1$}{
  \ForEach{client $k\!=\!1,\dots,K$ {in parallel}}{
    \tcc{{Step‑ahead start}}
    $w_{r+\frac12,0}^{(k)} \gets w_{r}^{(k)} - \alpha_r\,e_{r}^{(k)}$\;

    \tcc{\footnotesize {Local SGD with $T$ steps}}
    \For{$t \gets 0$ \KwTo $T-1$}{
      $w_{r+\frac12,t+1}^{(k)} \gets w_{r+\frac12,t}^{(k)} 
           - \eta_r\,\nabla f\!\bigl(w_{r+\frac12,t}^{(k)}; \zeta_{r,t}^{(k)}\bigr)$\;
    }
    \tcc{accumulated local update}
    $g_{r}^{(k)} \gets w_{r+\frac12,0}^{(k)} - w_{r+\frac12,T}^{(k)}$ \;
    $u_{r+1}^{(k)} \gets (1-\alpha_r)\,e_{r}^{(k)} + g_{r}^{(k)}$\;
    $e_{r+1}^{(k)} \gets u_{r+1}^{(k)} - \mathcal{C}\!\bigl(u_{r+1}^{(k)}\bigr)$\;
    {send} $\mathcal{C}\!\bigl(u_{r+1}^{(k)}\bigr)$ to server\;
  }

  \tcc{{Server‑side aggregation}}
  $u_{r+1} \gets \dfrac{1}{K}\,\sum_{k=1}^{K} \mathcal{C}\!\bigl(u_{r+1}^{(k)}\bigr)$\;
  \tcc{apply averaged update}
  $w_{r+1} \gets w_{r} - \eta\,u_{r+1}$\;
  {broadcast} $w_{r+1}$ to all clients\;
  \ForEach{client $k$}{ $w_{r+1}^{(k)} \gets w_{r+1}$\; }
}
\end{algorithm}}

In EF with \emph{biased} compression (e.g., Top-$k$), each client $k$ maintains a residual $e_r^{(k)}$ of \emph{unsent} coordinates. Although the stochastic gradient is \emph{computed} at the received global model $w_r$, the effective update is applied at a de-errored point $\tilde w_r := w_r - \delta_r$, where $\delta_r$ denotes the EF carry (often close to the mean residual $\bar e_r := \tfrac1K\sum_k e_r^{(k)}$ under biased sparsification). Thus, the gradient used for the step is \emph{stale} with respect to the point being updated, an effect akin to a \emph{staleness-of-one} delay in asynchronous SGD where gradients are evaluated at one iterate but applied to another \citep{lian2015asynchronous}. Formally, this mismatch is captured by \(g_k(\tilde w_r;\zeta) - g_k(w_r;\zeta)\), where \(g_k := \nabla f_k\), and under $L$-smoothness the local displacement error satisfies
\[
\varepsilon_r^{\mathrm{loc}}(0)
:= \frac{1}{K}\sum_{k=1}^K \mathbb{E}_\zeta\!\big\|g_k(\tilde w_r;\zeta)-g_k(w_r;\zeta)\big\|^2
\;\le\; L^2\|\delta_r\|^2 \;+\; 4\sigma^2.
\]

% In the classical gradient-compression setting with $T{=}1$, EF21 and its extensions~\citep{richtarik2021ef21,fatkhullin2025ef21bells} provide a modern error-feedback framework with clean contraction guarantees and no error floor, by compressing full gradients (or gradient differences) at a shared iterate and avoiding local drift. However, in our federated local-SGD regime, the compressed object is the accumulated local update over $T>1$ steps on heterogeneous data, and the model on each client drifts between communications. The EF update above coincides with the federated variant analyzed in Fed-EF~\citep{li2023analysis}, which we take as the EF-style backbone on top of which SA-PEF is built.

\paragraph{Algorithm overview.} 
SA-PEF (Algorithm~\ref{alg:sa_pef}) augments Local-SGD with two complementary ideas: 
(i) a \emph{step-ahead preview} of each client’s EF residual, and 
(ii) a \emph{partial carry-over} of that residual through standard EF. 
A per-round weight $\alpha_r \in [0,1]$ determines how much of the residual is previewed versus retained. In effect, SA-PEF \emph{previews a fraction of the residual and remembers the rest}, achieving fast early progress like SAEF while preserving the long-term stability of EF. Each client $k$ maintains a local model $w_r^{(k)}$ and a residual $e_r^{(k)}$, both initialized at round $r=0$.

At each communication round $r=0,\dots,R-1$, SA-PEF implements the following steps: 
\begin{enumerate}[leftmargin=*,itemsep=0pt,topsep=2pt,parsep=0pt]
  \item \emph{Step-ahead preview.} 
  Before local training, client $k$ shifts its model by a fraction $\alpha_r$ of its residual:
  \[
  w^{(k)}_{r+\tfrac12,0} \;=\; w^{(k)}_{r} - \alpha_r\, e^{(k)}_{r}.
  \]
  This moves gradient evaluation closer to where EF actually applies the update, improving
  alignment of the next direction with $-\nabla f(w_r)$.
  
  \item \emph{Local SGD.} 
  Starting from $w^{(k)}_{r+\tfrac12,0}$, client $k$ performs $T$ local SGD steps with stepsize $\eta_r$:
  \[
  w^{(k)}_{r+\tfrac12,t+1} \;=\; w^{(k)}_{r+\tfrac12,t}
  -\eta_r \nabla f\!\bigl(w^{(k)}_{r+\tfrac12,t};\zeta^{(k)}_{r,t}\bigr), 
  \quad t=0,\ldots,T-1.
  \]
  The accumulated local update is: $g_r^{(k)} \;=\; w^{(k)}_{r+\tfrac12,0} - w^{(k)}_{r+\tfrac12,T}.$
  
  \item \emph{Partial EF composition.} 
  SA-PEF blends the remaining residual with the new local update:
  \[
  u_{r+1}^{(k)} \;=\; (1-\alpha_r)\, e_r^{(k)} + g_r^{(k)}.
  \]
  The compressed message and residual update are then $\tilde u_{r+1}^{(k)} = \mathcal{C}\!\bigl(u_{r+1}^{(k)}\bigr), $
  $e_{r+1}^{(k)} = u_{r+1}^{(k)} - \tilde u_{r+1}^{(k)}.$
  Only $\tilde u_{r+1}^{(k)}$ is transmitted to the server and $e_{r+1}^{(k)}$ is retained locally.
  
  \item \emph{Server aggregation.} 
  The server averages compressed updates and applies them:
  \[
  u_{r+1} = \tfrac{1}{K}\sum_{k=1}^K \tilde u_{r+1}^{(k)}, 
  \qquad
  w_{r+1} = w_r - \eta \, u_{r+1}.
  \]
  The new global model $w_{r+1}$ is broadcast to all clients for the next round.
\end{enumerate}

\paragraph{Relation to EF and SAEF.} 
The step-ahead coefficient $\alpha_r$ allows SA-PEF to smoothly interpolate between prior methods:
\begin{itemize}[leftmargin=*,itemsep=0pt,topsep=0pt,parsep=0pt]
  \item $\alpha_r=0$: reduces to Fed-EF/classical EF in the federated local-SGD setting~\citep{li2023analysis}. 
  \item $\alpha_r=1$: reduces to a full step-ahead variant analogous to SAEF~\citep{xu2021step}. 
  \item $0<\alpha_r<1$: partial preview, fast early progress with reduced late-stage noise. 
\end{itemize}

A first-order view (descent lemma) gives $g_k(w_r-\alpha_r e_r^{(k)}) \approx g_k(w_r) - \alpha_r \nabla^2 f_k(w_r)\,e_r^{(k)}$, so $\mathbb{E}_r[\bar g_r(\alpha_r)] \approx \nabla f(w_r) - \alpha_r \bar H_r \bar e_r$ with $\bar H_r := \tfrac1K\sum_k \nabla^2 f_k(w_r)$. Thus preview provides a \emph{linear} alignment gain via $-\alpha_r \bar H_r \bar e_r$. On the other hand, the \emph{local-displacement} mismatch between $\nabla f_k(w_r)$ and $\nabla f_k(w_r-\alpha_r e_r^{(k)})$ grows \emph{quadratically} with $\alpha_r$ under $L$-smoothness. Hence, SA-PEF chooses an \emph{interior} $\alpha_r\in(0,1)$ to balance alignment benefits against the smoothness-driven cost, while preserving EF memory through $(1-\alpha_r)e_r^{(k)}$.

\paragraph{Residual contraction.} 
Our analysis yields a per-round residual contraction
\[
\rho_r=\Bigl(1-\tfrac{1}{\delta}\Bigr)\Bigl(2(1-\alpha_r)^2+24\,\alpha_r^2(\eta_r L T)^2\Bigr),
\]
which is strictly smaller than EF’s value $2(1-1/\delta)$ when $s_r=\eta_rLT$ is small, explaining SA-PEF’s faster early-phase convergence while retaining EF-style long-term stability.

% \paragraph{Practical considerations.} 
% Any $\delta$-contractive compressor (e.g., Top-$k$, quantization, sketching) is supported. SA-PEF requires no server-side changes beyond standard averaging, and clients only store a residual vector, making the method a \emph{lightweight drop-in replacement} for Local-SGD.


\section{Convergence Analysis}

\begin{definition}[Compression Operator]
\label{def:compression-operator}
A (possibly randomized) mapping $\mathcal{C}: \mathbb{R}^d \to \mathbb{R}^d$ is called a compression operator if there exists $\delta > 0$ such that, for all $w \in \mathbb{R}^d$,
\begin{equation}
\label{eq4}
\mathbb{E}\bigl[\|\mathcal{C}(w) - w\|_2^2\bigr]
\;\le\;
\Bigl(1 - \tfrac{1}{\delta}\Bigr)\,\|w\|_2^2,   \tag{4}
\end{equation}
where the expectation is taken over the internal randomness of $\mathcal{C}$.
\end{definition}

This class includes many commonly used compressors, such as Top-$k$, scaled Rand-$k$, and quantization operators. 
The parameter $\delta$ captures the contraction strength, with larger $\delta$ corresponding to weaker contraction (i.e., more aggressive compression).  

Our objective is to establish nonconvex convergence guarantees by upper‑bounding the average squared gradient norm:
\(\tfrac1R\sum_{r=0}^{R-1}\mathbb{E}\Vert\nabla f(w_r)\Vert^2\),
where expectations are taken over mini‑batch sampling, client selection, and compression randomness.

We adopt the standard assumption set used in \citet{alistarh2018convergence, li2023analysis, beznosikov2023biased} and follow the
notation from Alg.~\ref{alg:sa_pef}. Fix a communication round $r$ and let $\{\mathcal F_{r,t}\}_{t\ge 0}$ denote the natural filtration generated by all randomness up to the beginning of local step $t$ of round $r$.  At local step $t\in\{0,\dots,T-1\}$ on client $k$, the stochastic gradient takes the form
\(
g_{r,t}^{(k)} = \nabla f_k\left(w_{r+\frac12,t}^{(k)}\right) + \xi_{r,t}^{(k)},
\)
where $\xi_{r,t}^{(k)}$ captures stochastic noise.

\begin{assumption}[Smoothness]\label{A1}
Each local objective $f_k$ is differentiable with $L$-Lipschitz gradient: 
\(
\bigl\|\nabla f_k(y)-\nabla f_k(x)\bigr\|\le L\|y-x\|, \forall x,y\in\mathbb R^d.
\)
\end{assumption}

\begin{assumption}[Stochastic gradients]\label{A2}
For every client $k$ and step $t$, we have (i) $\mathbb E[\xi_{r,t}^{(k)}|\mathcal F_{r,t}]=0$ (unbiasedness), (ii) $\mathbb E[\|\xi_{r,t}^{(k)}\|^{2}|\mathcal F_{r,t}]\le \sigma^{2}$ (bounded variance), (iii) given $\mathcal F_{r,0}$, the family $\{\xi_{r,t}^{(k)}: k=1,\dots,K,\ t=0,\dots,T-1\}$ is conditionally independent.
\end{assumption}

\begin{assumption}[Gradient dissimilarity]\label{A3}
There exist constants $\beta^{2}\ge 1$ and $\nu^{2}\ge 0$ such that,
\[
\frac{1}{K}\sum_{k=1}^{K}\bigl\|\nabla f_k(x)\bigr\|^{2}
\;\le\;
\beta^{2}\,\bigl\|\nabla f(x)\bigr\|^{2} \;+\; \nu^{2},\qquad \forall x\in\mathbb R^{d}
\]
\end{assumption}

We now state a convergence guarantee for SA-PEF in our federated setting. 
%For clarity, the theorem is presented with \emph{constant} client stepsize $\eta_0$ and server stepsize $\eta$.
The full proof is provided in Appendix \ref{theorem_1_proof}.
\begin{theorem}[Stationary-point bound with constant inner-loop step]
\label{main_theorem}
Assume \ref{A1}--\ref{A3} and let the compressor satisfy Definition~\ref{def:compression-operator} with parameter $\delta\ge1$.
Run SA-PEF for $R\ge1$ rounds with a constant inner-loop stepsize $\eta_r\equiv\eta_0$, and set $s_0:=\eta_0LT\le \tfrac18$. Suppose further that
\(18\,\beta^{2}s_0^{2}\le \tfrac18\) and \(\eta\le \frac{1}{256\,\beta^{2}L\eta_0T}\).
Define the maximal residual contraction
\[
\rho_{\max}:=\sup_r\Bigl(1-\tfrac1\delta\Bigr)\Bigl(2(1-\alpha_r)^2+24\,\alpha_r^2 s_0^2\Bigr)<1,
\]
and the effective error constant
\[
\Theta:=\frac{16}{\eta}\times\frac{\mathcal E_{\max}}{1-\rho_{\max}}\Bigl(1-\tfrac1\delta\Bigr)\beta^{2}\,
\bigl(8\,\eta_0T+288\,L^{2}\eta_0^{3}T^{3}\bigr),
\]
where $\mathcal E_{\max}:=\sup_{r}\mathcal E_r$ is the maximum residual–error coefficient across rounds.
If $\Theta\le\tfrac12$, then with $f^\star:=\inf_x f(x)$ and initial residuals $e_0^{(k)}\equiv0$, we have
\begin{align*}
\frac{1}{R}\sum_{r=0}^{R-1}\E\|\nabla f(w_r)\|^{2}
\ \le\
&\frac{32\,(f(w_0)-f^\star)}{\eta\,\eta_0T\,R}
\;+\;\Bigl(1-\tfrac{1}{\delta}\Bigr)\Bigl[
C_\sigma\,\eta_0^{2}L^{2}T\,\sigma^{2}
+ C_\nu\,\eta_0^{2}L^{2}T^{2}\,\nu^{2}\Bigr]
\\[-2pt]
&\;+\;\frac{128\,L\,\eta_0}{K}\,\sigma^{2},
\end{align*}
where
\[
\begin{aligned}
C_\sigma&=\frac{32}{\eta}\Bigl[6\,\eta\,\eta_0^{2}L^{2}T+96\,L^{3}\eta\,\eta_0^{3}T^{2}\Bigr]
+\frac{32}{\eta}\times\frac{\mathcal E_{\max}}{1-\rho_{\max}}\Bigl[4\,\eta_0+96\,L^{2}\eta_0^{3}T^{2}\Bigr],\\
C_\nu&=\frac{32}{\eta}\Bigl[84\,\eta\,\eta_0^{2}L^{2}T^{2}+1344\,L^{3}\eta\,\eta_0^{3}T^{3}\Bigr]
+\frac{32}{\eta}\times\frac{\mathcal E_{\max}}{1-\rho_{\max}}\Bigl[8\,\eta_0T+1344\,L^{2}\eta_0^{3}T^{3}\Bigr].
\end{aligned}
\]
\end{theorem}

% \begin{remark}[Relation to EF21]
% EF21 and its extensions~\citep{richtarik2021ef21,fatkhullin2025ef21bells} obtain stronger guarantees (no error floor) under a different regime: synchronized data-parallel training with $T{=}1$, full-gradient (or gradient-difference) compression at a shared iterate, and no local steps. In this setting, EF21 is strictly preferable to classical EF. Our analysis targets a complementary regime, federated local-SGD with $T>1$ local steps, partial participation, and biased contractive compressors, where the compressed object is the accumulated local update and client drift plays a central role. Extending EF21-style arguments (or designing EF21-style step-ahead variants) to this local-SGD, partial-participation, biased-compression setting is non-trivial and remains an interesting direction for future work.
% \end{remark}

\paragraph{Discussion.}
Our result matches the standard nonconvex picture for \emph{biased} compression. With a constant inner stepsize $\eta_0$, the optimization error decreases at rate $O\big((\eta\,\eta_0 T R)^{-1}\big)$, while an $R$-independent floor remains due to residual drift. As in prior EF analyses, only the \emph{mini-batch variance} benefits from a $1/K$ reduction, where the residual-induced floor does not. Data heterogeneity contributes additively with a $T^2$ multiplier in the floor, while the stochastic variance floor carries a $T$ multiplier (both scaled by $\eta_0^2L^2$). The effect of compressor appears only via the usual bias factor $(1-1/\delta)$ and the residual contraction $\rho_{\max}<1$, which depends on $\alpha_r$ and $s_r=\eta_rLT$. This matches the qualitative dependence reported in earlier analyses of compressed FL \citep{li2023analysis, karimireddy2020scaffold}.

\begin{remark}[Partial participation (PP)] \label{rem:pp}
Let $p=m/K\in(0,1]$ be the participation rate, and assume constant client stepsize $\eta_r=\eta_0$ with $T$ local steps per round. Then, the averaged stationarity bound under partial participation is
\begin{align}
\frac{1}{R}\sum_{r=0}^{R-1}\E\|\nabla f(w_r)\|^{2}
\le
\mathcal{O}\bigg(\frac{f(w_0)-f_\star}{\eta\,p\,\eta_0 T\,R}\bigg)
+\mathcal{O}\bigg(\frac{L\,\eta_0}{\eta\,p\,m}\,\sigma^{2}\bigg) \notag
+\mathcal{O}\bigg(\frac{1}{\eta\,p}\Bigl(1-\tfrac{1}{\delta}\Bigr)\,
\eta_0^{2}L^{2}\bigl(T\,\sigma^{2}+T^{2}\nu^{2}\bigr)\bigg).
\end{align}
Thus, partial participation effectively reduces the horizon to $R_{\mathrm{eff}}=pR$: the optimization term slows by a factor $1/p$. The pure mini-batch variance averages down as $1/m$, whereas the compression/EF floors depend on $(1-1/\delta)$ and the local-work parameter $s=\eta_0LT$, and do not benefit from $1/m$ averaging.
\end{remark}

\paragraph{Comparison to EF under PP.}
With a diminishing stepsize chosen so that $\sum_{r=0}^{R-1}\eta_r T=\Theta(\sqrt{R})$, the optimization term scales as $\mathcal{O}\!\big(1/(p\sqrt{R})\big)$, or equivalently as $\mathcal{O}\!\big(\sqrt{K/m}/\sqrt{R}\big)$, which recovers the $\sqrt{K/m}$ slow-down of EF under partial participation \citep[Theorem~4.10]{li2023analysis}. Our variance terms likewise exhibit a $1/m$ reduction for mini-batch noise, while the compressor- and heterogeneity-induced floors remain of order $(1-1/\delta)$, in agreement with prior analyses.

\paragraph{Why step-ahead helps (and how much).}
Step-ahead modifies the residual contraction factor to
\[
\rho_r
=(1-\tfrac1\delta)\Bigl(2(1-\alpha_r)^2+24\,\alpha_r^2(\eta_rLT)^2\Bigr)
=(1-\tfrac1\delta)\Bigl(2-4\alpha_r+(2+24s_r^2)\alpha_r^2\Bigr),
\quad s_r:=\eta_rLT\le\tfrac18.
\]
This quadratic is minimized at
\(
\alpha_r^\star=\frac{1}{\,1+12s_r^2\,}\in(0.84,1],
\)
hence a \emph{moderate-to-large} step-ahead (close to $1$ when $s_r$ is small) achieves the strongest contraction. Relative to EF ($\alpha_r=0$), the contraction gain is
\(
\rho_r-\rho_{\mathrm{EF}}
=(1-\tfrac1\delta)[-4\alpha_r+(2+24s_r^2)\alpha_r^2]<0
\)
for $\alpha_r\in(0,\tfrac{1}{1+12s_r^2})$, with minimum value
\(
\rho_{\min}
=2(1-\tfrac1\delta)(1-\frac{1}{1+12s_r^2})
=\rho_{\mathrm{EF}}(1-\frac{1}{1+12s_r^2}).
\)
For $\alpha_r>\alpha_r^\star$, $\rho_r$ increases (since $\alpha_r^\star<1$), though the descent coupling in $(\eta,\eta_r)$ is unchanged. The trade-off is that larger $\alpha_r$ tightens the requirement on $s_r=\eta_rLT$ (via the $\alpha_r^2 s_r^2$ term) or necessitates milder compression (larger $\delta$) to maintain $\rho_r<1$. Overall, step-ahead reduces the contraction factor while leaving the leading optimization rate and the qualitative variance/heterogeneity terms unchanged.

\noindent\emph{Practical takeaway.}
When $s_r=\eta_rLT$ is small and $\alpha_r$ is chosen near $\alpha_r^\star$, the residual contraction factor is strictly smaller than in EF (cf.\ Prop.~\ref{prop:sapef-rho}), resulting in smaller constants in the residual–induced error terms. This leads to a steeper \emph{initial} decrease in the objective and gradient norm under the same stepsizes, which in turn yields faster convergence within a fixed communication budget. In regimes with high compression or strong data heterogeneity, SA-PEF can therefore outperform both standard EF and uncompressed Local-SGD. To place SA-PEF in context, we provide a brief comparison of Fed-EF, SAEF, CSER, SCAFCOM, and SA-PEF in Table~\ref{tab:rate-comparison}, highlighting that SA-PEF operates in the same FL regime as Fed-EF but achieves strictly better residual contraction under biased compression, while SCAFCOM relaxes heterogeneity assumptions at the cost of additional control-variates state.


\section{Experiments}

\subsection{Experimental setup}

We evaluate SA-PEF on three image classification benchmarks of increasing difficulty and scale. We use CIFAR-10~\citep{krizhevsky2009} with ResNet-9~\citep{page_cifar10_fast}, {CIFAR-100} \cite{krizhevsky2009} with {ResNet-18} \cite{he2016deep}, and {Tiny-ImageNet}~\citep{le2015tiny} with {ResNet-34}~\citep{he2016deep}, trained with cross-entropy loss. We apply standard preprocessing: per-dataset mean/std normalization. We create \(K=100\) clients and adopt \emph{partial participation} with rate \(p\in\{0.1,0.5,1.0\}\) where, in each round \(r\), the server samples \(m=\lfloor pK\rfloor\) clients uniformly without replacement. To induce client data heterogeneity, we apply \emph{Dirichlet} label partitioning with concentration parameter \(\gamma\in\{0.1,0.5,1.0\}\), where smaller \(\gamma\) indicates stronger non-IID~\citep{hsu2019measuring}. Each client’s local dataset remains fixed across rounds. Each selected client performs \(T=5\) local SGD steps per round. Training runs for \(R=200\) communication rounds. Unless otherwise stated, the local mini-batch size is 64, momentum is 0.9, and weight decay is \(5\times10^{-4}\) on CIFAR and \(1\times10^{-4}\) on Tiny-ImageNet. We use {Top-}\(k\) sparsification with sparsity level \(k/d\in\{0.01,0.05,0.1\}\). Clients transmit both indices and values of selected entries. We compare SA-PEF with {uncompressed LocalSGD}, {EF}~\citep{li2023analysis}, {SAEF}~\citep{xu2021step}, and {CSER}~\citep{xie2020cser}. All methods use the same client sampling, optimizer, learning-rate schedules, and total communication budget (rounds and bits). We report Top-1 accuracy versus rounds and communicated bits, rounds-to-target accuracy, and final accuracy at a fixed communication budget. We repeat all experiments with five random seeds and report mean values in the plots. For CIFAR-10 and CIFAR-100, we additionally report the final test accuracy as mean $\pm$ standard deviation over these five runs in Table~\ref{tab:error_bar}. Due to space constraints, we present results under high compression, with comprehensive results across all settings in the Appendix.

\subsection{Empirical results}

\captionsetup[subfigure]{justification=centering}

\begin{figure*}[t]
  \centering
  % Row 1
  \begin{subfigure}{0.245\textwidth}
    \includegraphics[width=\linewidth]{Plots/efficiency_plt_10_participation_top_1_resnet9_dir05.pdf}
    \caption{$p {=} 0.1$, $\gamma{=}0.5$}\label{fig:cifar10_10_5_e_main}
  \end{subfigure}\hfill
  \begin{subfigure}{0.245\textwidth}
    \includegraphics[width=\linewidth]{Plots/efficiency_plt_10_participation_top_10_resnet9_dir05.pdf}
    \caption{$p {=} 0.1$, $\gamma{=}0.5$}\label{fig:cifar10_10_5_e_m}
  \end{subfigure}\hfill
  \begin{subfigure}{0.245\textwidth}
    \includegraphics[width=\linewidth]{Plots/efficiency_plt_50_participation_top_1_resnet9.pdf}
    \caption{$p {=} 0.5$, $\gamma{=}0.1$}\label{fig:cifar10_50_1_e_main}
  \end{subfigure}
  \begin{subfigure}{0.245\textwidth}
    \includegraphics[width=\linewidth]{Plots/efficiency_plt_50_participation_top_10_resnet9.pdf}
    \caption{$p {=} 0.5$, $\gamma{=}0.1$}\label{fig:cifar10_50_10_e_main}
  \end{subfigure}
  \vspace{4pt}
  % Row 2
  \begin{subfigure}{0.245\textwidth}
    \includegraphics[width=\linewidth]{Plots/plt_10_participation_top_1_resnet9_dir05.pdf}
    \caption{$p {=} 0.1$, $\gamma{=}0.5$}\label{fig:cifar10_10_5_c_main}
  \end{subfigure}\hfill
  \begin{subfigure}{0.245\textwidth}
    \includegraphics[width=\linewidth]{Plots/plt_10_participation_top_10_resnet9_dir05.pdf}
    \caption{$p {=} 0.1$, $\gamma{=}0.5$}\label{fig:cifar10_10_5_c_m}
  \end{subfigure}\hfill
  \begin{subfigure}{0.245\textwidth}
    \includegraphics[width=\linewidth]{Plots/plt_50_participation_top_1_resnet9.pdf}
    \caption{$p {=} 0.5$, $\gamma{=}0.1$}\label{fig:cifar10_50_1_c}
  \end{subfigure}
  \begin{subfigure}{0.245\textwidth}
    \includegraphics[width=\linewidth]{Plots/plt_50_participation_top_10_resnet9.pdf}
    \caption{$p {=} 0.5$, $\gamma{=}0.1$}\label{fig:cifar10_50_10_c}
  \end{subfigure}
  \caption{Test accuracy vs number of rounds (row 1) and communicated GB (row 2) on the CIFAR-10 dataset using ResNet-9.}% across participation rates with Top-10\%
  \label{fig:cifar10_resnet9_main}
\end{figure*}

Figure~\ref{fig:cifar10_resnet9_main} compares SA-PEF with FedAvg (dense) and compressed baselines (EF, SAEF, CSER) on CIFAR-10 with ResNet-9 under two participation rates ($p\in\{0.1,0.5\}$) and two compression budgets (Top-1\%, Top-10\%).
In the accuracy versus rounds plots (top row), SA-PEF generally reaches a given accuracy in fewer rounds than EF and SAEF, with the largest margin in the harder regime ($\gamma{=}0.1$, Top-1\%). SAEF often shows an initial jump but tends to plateau, while SA-PEF continues to improve. CSER typically lags early and only catches up later.
In the accuracy versus communication plots (bottom row), SA-PEF’s curves are left-shifted: for the same test accuracy it requires less uplink communication (in Gigabyte) than EF or SAEF, whereas FedAvg attains high accuracy only at orders of magnitude higher cost.
Raising participation to $p{=}0.5$ benefits all approaches and narrows round-wise gaps, but SA-PEF remains the most communication-efficient across schemes.\footnote{Under low participation (e.g., 1-10\%), effective batch sizes shrink and both drift and compression noise increase, hence participation-aware hyperparameters (e.g., learning rate, local steps $T$, or $\alpha_r$) may need tuning. Here, we fix hyperparameters across methods for fairness, which can reduce the observable advantage of SA-PEF in extreme low-participation regimes.}

Figure~\ref{fig:cifar100_resnet18} shows results on CIFAR-100 with ResNet-18. Despite the increased task difficulty, the same qualitative trends persist: SA-PEF tends to dominate early rounds and delivers higher accuracy per unit of communication across most regimes, SAEF often plateaus early, and CSER improves mainly at larger communication budgets.\footnote{As in CIFAR-10, under low partial participation (1–10\%), the differences between methods may be less visible without participation-aware tuning of hyperparameters.}  The gains are most pronounced under aggressive compression (Top-1\%) and low participation (e.g., $p{=}0.1$).
Overall, these results suggest that combining preview with partial error feedback provides faster early progress and superior accuracy-communication trade-offs across architectures and datasets.
% comment from DKR: Under extreme partial participation (1-10\%), the effective batch size shrinks and client drift/compression noise grow thus, tuning of participation-aware hyperparameters (LR, T, $alpha_r$(may be)) are required, Otherwise the advantage of SA-PEF are largely washed out by sampling variance, and all method appear similarly.  I didn't get to tune these hyperparameters foe each setup as that requires a high compute. I can try to tune for SA-PEF but then it will be a bit unfair to other approaches. 

\begin{figure*}[t]
  \centering
  % Row 1
  \begin{subfigure}{0.245\textwidth}
    \includegraphics[width=\linewidth]{Plots/efficiency_plt_10_participation_top_1_resnet18.pdf}
    \caption{$p {=} 0.1$, Top-1\%}\label{fig:cifar100_10_1_e}
  \end{subfigure}\hfill
  \begin{subfigure}{0.245\textwidth}
    \includegraphics[width=\linewidth]{Plots/efficiency_plt_10_participation_top_10_resnet18.pdf}
    \caption{$p {=} 0.1$, Top-10\%}\label{fig:cifar100_10_10_e}
  \end{subfigure}\hfill
  \begin{subfigure}{0.245\textwidth}
    \includegraphics[width=\linewidth]{Plots//efficiency_plt_50_participation_top_1_resnet18.pdf}
    \caption{$p {=} 0.5$, Top-1\%}\label{fig:cifar100_50_1_e}
  \end{subfigure}
  \begin{subfigure}{0.245\textwidth}
    \includegraphics[width=\linewidth]{Plots/efficiency_plt_50_participation_top_10_resnet18.pdf}
    \caption{$p {=} 0.5$, Top-10\%}\label{fig:cifar100_50_10_e}
  \end{subfigure}
  \vspace{4pt}
  % Row 2
  \begin{subfigure}{0.245\textwidth}
    \includegraphics[width=\linewidth]{Plots/plt_10_participation_top_1_resnet18.pdf}
    \caption{$p {=} 0.1$, Top-1\%}\label{fig:cifar100_10_1_c}
  \end{subfigure}\hfill
  \begin{subfigure}{0.245\textwidth}
    \includegraphics[width=\linewidth]{Plots/plt_10_participation_top_10_resnet18.pdf}
    \caption{$p {=} 0.1$, Top-10\%}\label{fig:cifar100_10_10_c}
  \end{subfigure}\hfill
  \begin{subfigure}{0.245\textwidth}
    \includegraphics[width=\linewidth]{Plots/plt_50_participation_top_1_resnet18.pdf}
    \caption{$p {=} 0.5$, Top-1\%}\label{fig:cifar100_50_1_c}
  \end{subfigure}
  \begin{subfigure}{0.245\textwidth}
    \includegraphics[width=\linewidth]{Plots/plt_50_participation_top_10_resnet18.pdf}
    \caption{$p {=} 0.5$, Top-10\%}\label{fig:cifar100_50_10_c}
  \end{subfigure}
  \caption{Test accuracy vs number of rounds (row 1) and communicated GB (row 2) on the CIFAR-100 dataset using ResNet-18 and $\gamma{=}0.1$.}% across participation rates with Top-10\%
  \label{fig:cifar100_resnet18}
\end{figure*}

\paragraph{Discussion: } Overall, our results position SA-PEF as a lightweight but effective upgrade of classical EF in FL. Compared to EF and its step-ahead variant SAEF, SA-PEF converges consistently faster and offers better accuracy-communication trade-offs under practical settings. Relative to CSER, which periodically resets residuals to control mismatch, SA-PEF achieves comparable or better robustness without introducing any additional reset-period hyperparameter, and it avoids CSER’s high \emph{peak} communication cost when compressed or full residuals are transmitted at reset rounds. Since practical systems must provision for peak, rather than average, bandwidth and latency, this makes SA-PEF more attractive as a drop-in component in resource-constrained deployments. Control-variate methods such as SCAFFOLD and SCAFCOM target a complementary axis, mitigating client drift and heterogeneity via additional per-client state, whereas SA-PEF focuses on reducing compression-induced residual mismatch within the EF family. In this sense, SA-PEF is best viewed as a drop-in improvement for EF-style compressed FL (and a natural building block for future combinations with control variates), providing significant gains in challenging regimes with high compression and strong heterogeneity.

\subsection{Sensitivity Analysis of Step-Ahead Coefficient $\alpha$}
To assess the robustness of SA-PEF to the choice of step-ahead coefficient, we sweep $\alpha_r$ between zero and one with increment of $0.1$ on CIFAR-10 with ResNet-9 and CIFAR-100 with ResNet-18 under non-IID Dirichlet partitioning ($\gamma=0.1$), Top-1\% sparsification, $T=5$ local steps, and $p=0.1$ participation. Figure~\ref{alpha_ablation} reports test accuracy versus rounds for EF ($\alpha_r=0$), SAEF ($\alpha_r=1$), and SA-PEF with intermediate $\alpha_r$ values. Three regimes emerge: (i) \emph{Small $\alpha$} (\(\alpha_r \le 0.3\)) behaves similarly to EF, with noticeably slower convergence and lower final accuracy. (ii) \emph{Intermediate $\alpha$} (\(\alpha_r \in [0.6,0.9]\)) produces nearly identical curves, yielding the fastest convergence and highest final accuracy. This interval includes the default $\alpha_r=0.85$ used in our main experiments. (iii) \emph{Full step-ahead} (\(\alpha_r = 1.0\), SAEF) accelerates early rounds but plateaus slightly below the best SA-PEF setting in the later phases. Overall, SA-PEF is robust across a broad high-$\alpha$ region, while performance is significantly affected only at the extremes: $\alpha_r \approx 0$ (reducing to EF) or $\alpha_r = 1$ (SAEF). This supports treating $\alpha$ as a momentum-like parameter and use a single default value (e.g., $\alpha_r \approx 0.8$-$0.9$) across tasks without heavy tuning.

\begin{figure*}[h]
  \centering
  \begin{subfigure}{0.45\textwidth}
    \includegraphics[width=\linewidth]{Plots/alpha_ablation.pdf}
  \caption{with CIFAR-10 using ResNet-9}\label{fig:cifar10_10_1_alpha}
\end{subfigure}\hfill
\begin{subfigure}{0.45\textwidth}
  \includegraphics[width=\linewidth]{Plots/alpha_ablation_resnet18.pdf}
  \caption{with CIFAR-100 using ResNet-18}\label{fig:cifar100_10_1_alpha}
\end{subfigure}
\caption{Sensitivity analysis of step-ahead coefficient $\alpha$.}
\label{alpha_ablation}
\end{figure*}

\subsection{Gradient Mismatch}

Let $w\in\mathbb{R}^d$ stack all trainable parameters, and let $f(\cdot;w)$ be the network in \emph{evaluation} mode (dropout disabled, BatchNorm with frozen statistics). For a mini-batch $S=\{(x_i,y_i)\}_{i=1}^b$ drawn from a held-out loader, define the batch loss as
\(
\mathcal{L}_S(w)=\tfrac{1}{b}\sum_{(x,y)\in S}\ell\big(f(x;w),y\big).
\)
At round $r$, for client $k$ and $\alpha\in[0,1]$, we consider two evaluation points \(w_r\) and \(w_r^{(\alpha,k)}= w_r - \alpha\, e_r^{(k)}\).
Using the same mini-batch $S$ for both, we compute the associated gradients as
\(g_r= \nabla_w \mathcal{L}_S(w_r)\) and \(g_r^{(\alpha,k)}= \nabla_w \mathcal{L}_S\big(w_r^{(\alpha,k)}\big)\),
and define the squared gradient-mismatch as
\(
\hat\varepsilon_r^{(k)}(\alpha)=\big\|\, g_r - g_r^{(\alpha,k)}\big\|_2^{2}
\)
and its client average as $\hat\varepsilon_r(\alpha)=\tfrac{1}{K}\sum_{k=1}^K \hat\varepsilon_r^{(k)}(\alpha)$. In particular, we switch the model to \texttt{eval} mode, clear any stale gradients, reuse the same $S$, compute $g_r$ and $g_r^{(\alpha,k)}$ via first-order autodiff, and then restore the model to training mode. We do not alter any optimizer state or BatchNorm buffer.

% \begin{figure}[h]
% \begin{center}
% \includesvg[width=.6\columnwidth]{Plots/grad_mismatch_plot}
% \end{center}
% \caption{Gradient mismatch for different algorithms on the CIFAR-10 dataset using ResNet-9}
% \label{fig:mismatch}
% \end{figure}

\begin{figure*}[h]
  \centering
  % Row 1
  \begin{subfigure}{0.45\textwidth}
    \includegraphics[width=\linewidth]{Plots/grad_mismatch_plot_1.pdf}
    \caption{with CIFAR-10 using ResNet-9}\label{fig:cifar10_10_1_grad}
  \end{subfigure}\hfill
  \begin{subfigure}{0.45\textwidth}
    \includegraphics[width=\linewidth]{Plots/grad_mismatch_plot_2.pdf}
    \caption{with CIFAR-100 using ResNet-18}\label{fig:cifar100_10_1_grad}
  \end{subfigure}
  \caption{Gradient mismatch for different algorithms.}% across participation rates with Top-10\%
  \label{fig:mismatch}
\end{figure*}

Figure~\ref{fig:mismatch} shows the gradient-mismatch probe across training rounds for all methods. SA-PEF keeps the mismatch essentially flat and near zero throughout training: previewing only a fraction of the residual shifts the evaluation point closer to where EF applies the update, while the retained $(1-\alpha)e_r$ further prevents residual build-up. EF exhibits a steady late-phase rise, consistent with its one-step staleness effect as displacement accumulates over rounds. SAEF produces the largest spikes and highest overall mismatch, as with full preview ($\alpha{=}1$), the evaluation occurs at $w_r-e_r$. Large or heterogeneous residuals spike mismatch, causing late-stage plateaus. CSER stays near zero, reflecting its error-reset/averaging mechanism that suppresses persistent residual drift. Overall, the results confirm that combining partial preview with partial EF controls mismatch throughout training, unlike plain EF or SAEF.


\section{Conclusion}

We presented {SA-PEF}, which combines a step-ahead preview with partial error feedback to accelerate early-round progress in FL while preserving the long-run stability of EF under biased compression. Our theoretical analysis covers non-convex objectives, local SGD, partial participation, and \(\delta\)-contractive compressors, and establishes convergence to stationarity with rates matching non-convex FedSGD up to constants. To our knowledge, this is the first analysis of SAEF in a federated setting that simultaneously accounts for local updates, non-IID data, partial participation, and biased compression. Empirically, across datasets, models, and compressors, SA-PEF consistently reaches target accuracy in fewer rounds than EF and avoids the late-stage plateaus observed with full step-ahead. Looking ahead, a promising direction is developing \emph{adaptive schedules for \(\alpha\)}, guided by online estimates of noise, residual norms, or client drift, and analyzing the feedback between mismatch reduction and residual dynamics under partial participation. Another direction is integrating SA-PEF with \emph{adaptive optimizers and momentum}, including a preconditioned drift analysis to control momentum-residual interactions. Finally, a natural direction for future work is to bring ideas from EF21 into the federated local-SGD regime we study. %~\citep{richtarik2021ef21,fatkhullin2025ef21bells}
%e.g., by designing and analyzing EF21-style step-ahead variants for $T>1$ local steps, partial participation, and biased compression. 
We expect these extensions to further improve early-phase alignment while retaining stability at scale, advancing communication-efficient FL under realistic constraints. 


\paragraph{Ethics statement.}
This work studies communication-efficient federated optimization using public, non-sensitive benchmark datasets (CIFAR-10/100 and Tiny-ImageNet). No personal, identifiable, or protected-class attributes are used, and no user-generated private data was accessed. All datasets are widely used for research and were obtained under their respective licenses; we cite the original sources. Our algorithms are intended to \emph{reduce} communication and potentially lower the energy costs of federated training.


\paragraph{Reproducibility statement.}
We provide a zip file with anonymized code and configuration files. The package includes:
(i) exact configs for all experiments;
(ii) scripts to download datasets and to create federated partitions (IID and Dirichlet with $\alpha\in\{0.1,0.5\}$);
(iii) model definitions (ResNet-9/18/34), compressor implementations (Top-$k$ with $k/d\in\{1\%,10\%\}$, error-feedback variants), and our SA-PEF scheduler.
Unless otherwise stated, experiments use client participation $q\in\{0.1,0.5\}$, local epochs $T\in\{1,\ldots,5\}$, batch size $B=64$, optimizer \texttt{SGD} (momentum $0.9$), weight decay $5\!\times\!10^{-4}$, and base step size $\eta_0=0.1$ with cosine or step decay. Communication is measured as uplink GB, including indices and values for sparse updates. Hardware: runs were executed on A100/A5000/H200 GPUs.

% \begin{figure}[h]
% \begin{center}
% %\framebox[4.0in]{$\;$}
% \fbox{\rule[-.5cm]{0cm}{4cm} \rule[-.5cm]{4cm}{0cm}}
% \end{center}
% \caption{Sample figure caption.}
% \end{figure}

% The table number and title always appear before the table. See
% Table~\ref{sample-table}.

% \begin{table}[t]
% \caption{Sample table title}
% \label{sample-table}
% \begin{center}
% \begin{tabular}{ll}
% \multicolumn{1}{c}{\bf PART}  &\multicolumn{1}{c}{\bf DESCRIPTION}
% \\ \hline \\
% Dendrite         &Input terminal \\
% Axon             &Output terminal \\
% Soma             &Cell body (contains cell nucleus) \\
% \end{tabular}
% \end{center}
% \end{table}

%\subsubsection*{Author Contributions}
%If you'd like to, you may include  a section for author contributions as is done
%in many journals. This is optional and at the discretion of the authors.

%\subsubsection*{Acknowledgments}
%Use unnumbered third level headings for the acknowledgments. All
%acknowledgments, including those to funding agencies, go at the end of the paper.


\bibliography{iclr2026_conference}
\bibliographystyle{iclr2026_conference}

\newpage
\appendix
\section{Proof of Convergence Results}

\subsection{Proof of Theorem~\ref{main_theorem}}
\label{theorem_1_proof}
\begin{proof}
Subtracting the \emph{average residual} from the communicated model converts a biased, compressed update into something that resembles vanilla SGD.  

We have:
\begin{enumerate}
  \item \textbf{Step-ahead shift} $w_{r+\frac12}^{(k)}\!=w_r^{(k)}-\alpha_r e_r^{(k)}$.
  \item \textbf{Residual update \& compression} $e_{r+1}^{(k)}=u_{r+1}^{(k)}-C(u_{r+1}^{(k)})$ with\\
  $u_{r+1}^{(k)}=(1-\alpha_r)e_r^{(k)}+g_r^{(k)}$.
\end{enumerate}

Let
\[
\bar e_r:=\frac1K\sum_{k=1}^Ke_r^{(k)},\quad
\bar u_{r+1}:=(1-\alpha_r)\bar e_r+\bar g_r,\quad
\bar C_{r+1}:=\frac1K\sum_{k}C(u_{r+1}^{(k)}) ,
\]
so the server update is $w_{r+1}=w_r-\eta\;\bar C_{r+1}$.

Define a virtual iterate:
\[
\begin{aligned}
x_r
  &= w_r-\eta\,\bar e_r \\
x_{r+1}
  &= w_{r+1}-\eta\,\bar e_{r+1}                                                   \notag\\
  &= (w_r-\eta\,\bar C_{r+1})-\eta\,\bigl(\bar u_{r+1}-\bar C_{r+1}\bigr)               \\
  &= w_r-\eta\,\bar u_{r+1}                                                       \\
  &= (x_r+\eta\,\bar e_r)\;-\;\eta\,\bigl[(1-\alpha_r)\bar e_r+\bar g_r\bigr]           \\
  &= x_r\;+\;\eta\,\alpha_r\bar e_r\;-\;\eta\,\bar g_r. 
\end{aligned} 
\]

\begin{equation}\label{eq:descent_raw}
   f(x_{r+1}) \;\le\; f(x_r)
      + \bigl\langle \nabla f(x_r),\,x_{r+1}-x_r \bigr\rangle
      + \frac{L}{2}\,\bigl\|x_{r+1}-x_r\bigr\|^{2}.
\end{equation}

Because \(x_{r+1}-x_r = \eta\bigl(\alpha_r\bar e_r-\bar g_r\bigr)\), taking the
expectation over all randomness inside round \(r\) yields
\begin{align}\label{eq:descent_exp}
   \mathbb{E}\!\bigl[f(x_{r+1})\bigr]
   \;\le\;
   \mathbb{E}\!\bigl[f(x_r)\bigr]
   + \eta\,\mathbb{E}\!\bigl[
       \bigl\langle\nabla f(x_r),\,\alpha_r\bar e_r-\bar g_r\bigr\rangle
     \bigr]
   + \frac{\eta^{2}L}{2}\,
     \mathbb{E}\!\bigl[
       \bigl\|\alpha_r\bar e_r-\bar g_r\bigr\|^{2}
     \bigr].
\end{align}

Isolating the drift of the objective between two consecutive rounds and
decomposing the inner product gives
\begin{align}\label{eq:descent_decomp}
   \mathbb{E}\!\bigl[f(x_{r+1})\bigr]
   - \mathbb{E}\!\bigl[f(x_r)\bigr]
   &\le
      \eta\,\mathbb{E}\!\bigl[
        \bigl\langle\nabla f(x_r),\,\alpha_r\bar e_r-\bar g_r\bigr\rangle
      \bigr]
      + \frac{\eta^{2}L}{2}\,
        \mathbb{E}\!\bigl[
          \bigl\|\alpha_r\bar e_r-\bar g_r\bigr\|^{2}
        \bigr]                                         \notag   \\[2pt]
   &=
      \underbrace{\eta\,\mathbb{E}\!\bigl[
        \bigl\langle\nabla f(w_r),\,\alpha_r\bar e_r-\bar g_r\bigr\rangle
      \bigr]}_{T1} 
      + \underbrace{\frac{\eta^{2}L}{2}\,
        \mathbb{E}\!\bigl[
          \bigl\|\alpha_r\bar e_r-\bar g_r\bigr\|^{2}
        \bigr]}_{T2}                                        \notag     \\[2pt]
   &\quad
      + \underbrace{\eta\,\mathbb{E}\!\bigl[
        \bigl\langle
          \nabla f(x_r)-\nabla f(w_r),\,\alpha_r\bar e_r-\bar g_r
        \bigr\rangle
      \bigr]}_{T3}.
\end{align}
\paragraph{Bounding the first inner-product term}
Recall
\[
T_{1}
=\eta\,\E_{r}\!\bigl[\langle\nabla f(w_{r}),\,\alpha_{r}\bar e_{r}-\bar g_{r}\rangle\bigr]
= \eta\,\alpha_r\,\langle\nabla f(w_r),\bar e_r\rangle
   -\eta\,\E_r\!\bigl[\langle\nabla f(w_r),\bar g_r\rangle\bigr],
\]
where
\(
\bar g_r=\frac{\eta_r}{K}\sum_{k=1}^K\sum_{t=0}^{T-1}(\nabla f_k(w_{r+\frac12,t}^{(k)})+\xi_{k,t})
\)
and \(\E_r[\xi_{k,t}]=0\).\\
Define \(\bar\nabla_t := \frac1K\sum_{k=1}^K\nabla f_k(w_{r+\frac12,t}^{(k)})\) and \(\Delta_t = \frac1K\sum_{k=1}^K(\nabla f_k(w_{r+\frac12,t}^{(k)})-\nabla f(w_r))\) \\
Write \(\bar\nabla_t=\nabla f(w_r)+\Delta_t\) and use \(L\)-smoothness to get
\(
\E_r\|\Delta_t\|^2
\le L^2\times \frac1K\sum_k\E_r\|w_{r+\frac12,t}^{(k)}-w_r\|^2.
\) \\
Young’s inequality gives, for any \(\lambda_1>0\),
\[
\E_r\langle\nabla f(w_r),\bar g_r\rangle
\ge \eta_r T\Bigl(1-\frac{\lambda_1}{2}\Bigr)\|\nabla f(w_r)\|^2
-\frac{\eta_r L^2}{2\lambda_1}\,S,
\quad
S:=\frac1K\sum_{k=1}^K\sum_{t=0}^{T-1}\E_r\|w_{r+\frac12,t}^{(k)}-w_r\|^2.
\]
By Lemma~\ref{lem:local-drift} (summed over \(t\)),
\[
S \le
12\,\eta_r^{2}T^{2}\sigma^{2}
+168\,\eta_r^{2}T^{3}\nu^{2}
+36\,\eta_r^{2}T^{3}\beta^{2}\|\nabla f(w_r)\|^{2}
+3T\alpha_r^2 \bar E_r.
\]
For the step-ahead piece, Young’s inequality yields, for any \(\lambda_2>0\),
\[
\eta\,\alpha_r\,\langle\nabla f(w_r),\bar e_r\rangle
\le \frac{\eta\lambda_2}{2}\|\nabla f(w_r)\|^2
   +\frac{\eta\,\alpha_r^2}{2\lambda_2}\,\bar E_r.
\]
Combining,
\[
T_1
\le
\eta\Bigl(\frac{\lambda_2}{2}-\eta_r T\Bigl(1-\frac{\lambda_1}{2}\Bigr)\Bigr)\|\nabla f(w_r)\|^2
+\frac{\eta\,\alpha_r^2}{2\lambda_2}\,\bar E_r
+\frac{\eta\,\eta_r L^2}{2\lambda_1}\,S.
\]
Choose \(\lambda_1=1\) and \(\lambda_2=\eta_r T/2\); then
\[
T_1
\le -\frac{\eta\,\eta_r T}{4}\,\|\nabla f(w_r)\|^2
   +\frac{\eta\,\alpha_r^2}{\eta_r T}\,\bar E_r
   +\frac{\eta\,\eta_r L^2}{2}\,S,
\]
and substituting the bound on \(S\) gives
\begin{align}
T_1
&\le \eta\,\eta_r\Bigl[
   -\frac{T}{4}
   + 18\,\eta_r^{2}L^{2}\beta^{2}T^{3}
\Bigr]\,
\|\nabla f(w_r)\|^{2} \notag\\
&\qquad
+ \eta\,\alpha_r^{2}\Bigl[\frac{1}{\eta_r T}+\frac{3}{2}\eta_r L^{2}T\Bigr]\bar E_r
+ 6\,\eta\,\eta_r^{3}L^{2}T^{2}\,\sigma^{2}
+ 84\,\eta\,\eta_r^{3}L^{2}T^{3}\,\nu^{2}.
\label{eq:T1-final-self}
\end{align}



\textbf{Bounding the quadratic (smoothness) term}

Recall the second contribution in the descent inequality
\begin{equation}
T_{2}\;=\;\frac{L\eta^{2}}{2}\,
           \mathbb{E}_{r}\!\bigl[\,
               \|\alpha_{r}\bar e_{r}-\bar g_{r}\|^{2}
           \bigr].
\tag{8}
\label{P-T2}
\end{equation}

\paragraph{Bounding the gradient–mismatch term}
\[
T_{3}
:=\;
\eta\,
\E_{r}\!\bigl[
  \langle\nabla f(x_{r})-\nabla f(w_{r}),\,
          \alpha_{r}\bar e_{r}-\bar g_{r}\rangle
\bigr].
\]
Since $x_r = w_r - \eta \bar e_r$, we have $\|x_r-w_r\|=\eta\|\bar e_r\|$, and by $L$-smoothness,
\[
\|\nabla f(x_r)-\nabla f(w_r)\|\le L\|x_r-w_r\|=L\eta\|\bar e_r\|.
\]
Apply Young’s inequality (any $\lambda>0$) with
\(a=\nabla f(x_r)-\nabla f(w_r)\) and \(b=\alpha_r\bar e_r-\bar g_r\):
\[
\bigl|\langle\nabla f(x_r)-\nabla f(w_r),\,\alpha_r\bar e_r-\bar g_r\rangle\bigr|
\le
\frac{\lambda}{2}\,L^{2}\eta^{2}\,\|\bar e_r\|^2
+\frac{1}{2\lambda}\,\|\alpha_r\bar e_r-\bar g_r\|^2.
\]
Taking $\E_r[\cdot]$, using $\|\bar e_r\|^2\le \bar E_r$, and multiplying by $\eta$ yields
\[
|T_3|
\;\le\;
\frac{\lambda L^{2}\eta^{3}}{2}\,\bar E_r
\;+\;
\frac{\eta}{2\lambda}\,\E_r\!\bigl[\|\alpha_r\bar e_r-\bar g_r\|^2\bigr].
\]
Choosing $\displaystyle \lambda=\frac{1}{L\eta}$ balances the two terms and $T_3\le |T_3|$ gives

\begin{align}
T_3
&\le \frac{L\eta^2}{2}\,\bar E_r
   + \frac{L\eta^2}{2}\,\E_r\!\bigl[\|\alpha_r\bar e_r-\bar g_r\|^2\bigr].
\label{eq:T3}
\end{align}

Put the pieces together:

\begin{align}
\E_r\!\bigl[f(x_{r+1})-f(x_r)\bigr]
&\le \eta\,\eta_r\Bigl[
   -\frac{T}{4}
   + 18\,\eta_r^{2}L^{2}\beta^{2}T^{3}
\Bigr]\,
\|\nabla f(w_r)\|^{2} \notag\\
&\qquad
+ \eta\,\alpha_r^{2}\Bigl[\frac{1}{\eta_r T}+\frac{3}{2}\eta_r L^{2}T\Bigr]\bar E_r
+ 6\,\eta\,\eta_r^{3}L^{2}T^{2}\,\sigma^{2}
+ 84\,\eta\,\eta_r^{3}L^{2}T^{3}\,\nu^{2} \notag\\
&\qquad +\;L\eta^{2}\,
           \mathbb{E}_{r}\!\bigl[\,
               \|\alpha_{r}\bar e_{r}-\bar g_{r}\|^{2}
           \bigr] + \frac{L\eta^2}{2}\,\bar E_r
\end{align}

Using Lemma~\ref{lem:step-ahead-moment-fixed},
\[
\begin{aligned}
\E_r\!\bigl[\|\alpha_r\bar e_r-\bar g_r\|^{2}\bigr]
&\le 
2\alpha_r^{2}\bar E_r
+ 8\eta_r^{2}T^{2}\|\nabla f(w_r)\|^{2}\bigl(1+36L^{2}\eta_r^{2}T^{2}\beta^{2}\bigr) \\
&\qquad
+\;\frac{4\eta_r^{2}T}{K}\,\sigma^{2}
+\;96L^{2}\eta_r^{4}T^{3}\bigl(\sigma^{2}+14T\nu^{2}\bigr)
+\;24\alpha_r^{2}\eta_r^{2}L^{2}T^{2}\bar E_r.
\end{aligned}
\]

\begin{align}
\E_r\!\bigl[f(x_{r+1})-f(x_r)\bigr]
&\le \Bigl[\eta\,\eta_r\!\Bigl(-\tfrac{T}{4}
+ 18\,\eta_r^{2}L^{2}\beta^{2}T^{3}\Bigr)
\;+\; L\eta^{2}\Bigl(8\eta_r^{2}T^{2}
+ 288\,L^{2}\eta_r^{4}T^{4}\beta^{2}\Bigr)\Bigr]\|\nabla f(w_r)\|^{2}\notag\\[2pt]
&+
\Bigl[\eta\,\alpha_r^{2}\Bigl(\tfrac{1}{\eta_r T}+\tfrac{3}{2}\eta_r L^{2}T\Bigr)
\;+\; L\eta^{2}\Bigl(2\alpha_r^{2}
+24\,\alpha_r^{2}\eta_r^{2}L^{2}T^{2}\Bigr)
\;+\;\tfrac{L\eta^{2}}{2}\Bigr]\bar E_r, \notag\\[2pt]
&+\Bigl[
6\,\eta\,\eta_r^{3}L^{2}T^{2}
+ \tfrac{4L\eta^{2}\eta_r^{2}T}{K}
+ 96\,L^{3}\eta^{2}\eta_r^{4}T^{3}
\Bigr]\,\sigma^{2} \notag\\[2pt]
&+\Bigl[
84\,\eta\,\eta_r^{3}L^{2}T^{3}
+ 1344\,L^{3}\eta^{2}\eta_r^{4}T^{4}
\Bigr]\,\nu^{2}
\end{align}


\paragraph{Telescoping.}
Summing over $r=0,\dots,R-1$ and taking total expectation,
\begin{align}
\E\!\bigl[f(x_{R})-f(x_0)\bigr]
&\le \sum_{r=0}^{R-1}\Bigl\{
A_r\,\E\|\nabla f(w_r)\|^{2}
+ E_r\,\E[\bar E_r]
+ V_r\,\sigma^{2}
+ H_r\,\nu^{2}\Bigr\},\label{eq:telescoped-raw}
\end{align}
where
\[
\begin{aligned}
A_r
&:= \eta\,\eta_r\!\Bigl(-\tfrac{T}{4}
+ 18\,\eta_r^{2}L^{2}\beta^{2}T^{3}\Bigr)
+ L\eta^{2}\!\Bigl(8\eta_r^{2}T^{2}
+ 288\,L^{2}\eta_r^{4}T^{4}\beta^{2}\Bigr),\\[2pt]
E_r
&:= \eta\,\alpha_r^{2}\Bigl(\tfrac{1}{\eta_r T}+\tfrac{3}{2}\eta_r L^{2}T\Bigr)
+ L\eta^{2}\Bigl(2\alpha_r^{2}
+24\,\alpha_r^{2}\eta_r^{2}L^{2}T^{2}\Bigr)
+ \tfrac{L\eta^{2}}{2},\\[2pt]
V_r
&:= 6\,\eta\,\eta_r^{3}L^{2}T^{2}
+ \tfrac{4L\eta^{2}\eta_r^{2}T}{K}
+ 96\,L^{3}\eta^{2}\eta_r^{4}T^{3},\\[2pt]
H_r
&:= 84\,\eta\,\eta_r^{3}L^{2}T^{3}
+ 1344\,L^{3}\eta^{2}\eta_r^{4}T^{4}.
\end{aligned}
\]

We now derive the main convergence guarantee. The analysis begins with the telescoped recursion from \eqref{eq:telescoped-raw}, which provides an upper bound on the function value progress, $\E[f(x_R) - f(x_0)]$. The key step is to ensure that the coefficient $A_r$ of the squared gradient norm is negative, which guarantees descent. The following lemma establishes sufficient conditions for this.

\begin{lemma}[Sufficient Conditions for Descent]
\label{lem:descent_conditions}
Let $s_r := \eta_r L T$. Assume that for all rounds $r$, the parameters satisfy $s_r \le 1/8$ and the following two conditions:
\begin{align}
    18\,\beta^{2}\,s_r^{2} &\le \frac{1}{8}, \label{eq:cond1} \\
    \eta &\le \frac{1}{256\,\beta^{2}\,L\,\eta_r T}, \label{eq:cond2}
\end{align}
where it is assumed that $\beta^2 \ge 1$. Then, the coefficient $A_r$ is bounded as
\[
    A_r \le -\frac{1}{16}\,\eta\,\eta_r T.
\]
\end{lemma}

\begin{proof}
The coefficient $A_r$ is defined as
\[
A_r = -\frac{\eta\eta_r T}{4} + 18\eta\eta_r^3 L^2 \beta^2 T^3 + L\eta^2(8\eta_r^2 T^2 + 288 L^2 \eta_r^4 T^4 \beta^2).
\]
We proceed by bounding the two positive terms separately.

First, we bound the term $18\eta\eta_r^3 L^2 \beta^2 T^3$. By applying the definition $s_r = \eta_r L T$ and condition \eqref{eq:cond1}, we obtain
\begin{align*}
    18\eta\eta_r^3 L^2 \beta^2 T^3 &= (\eta\eta_r T) \times (18 \beta^2 \eta_r^2 L^2 T^2) \\
    &= (\eta\eta_r T) \times (18 \beta^2 s_r^2) \le \frac{1}{8}\eta\eta_r T.
\end{align*}

Next, we bound the term $L\eta^2(8\eta_r^2 T^2 + 288 L^2 \eta_r^4 T^4 \beta^2)$. Condition \eqref{eq:cond1} implies $36\beta^2 s_r^2 \le 1/4$. This allows us to simplify the expression within parentheses:
\begin{align*}
    8\eta_r^2 T^2 + 288 L^2 \eta_r^4 T^4 \beta^2 &= 8\eta_r^2 T^2 \left(1 + 36 \beta^2 L^2 \eta_r^2 T^2 \right) \\
    &= 8\eta_r^2 T^2 (1 + 36 \beta^2 s_r^2) \\
    &\le 8\eta_r^2 T^2 \left(1 + \frac{1}{4}\right) = 10\eta_r^2 T^2.
\end{align*}
Using this intermediate result and condition \eqref{eq:cond2}, we bound the full term. The assumption $\beta^2 \ge 1$ ensures that $256\beta^2 \ge 160$.
\begin{align*}
    L\eta^2(8\eta_r^2 T^2 + 288 L^2 \eta_r^4 T^4 \beta^2) &\le L\eta^2(10\eta_r^2 T^2) \\
    &= (10 L \eta \eta_r^2 T^2) \times \eta \\
    &\le 10 L \eta \eta_r^2 T^2 \times \left( \frac{1}{256\beta^2 L \eta_r T} \right) \\
    &= \frac{10}{256\beta^2} \eta \eta_r T \le \frac{10}{160} \eta \eta_r T = \frac{1}{16} \eta \eta_r T.
\end{align*}

Finally, substituting these bounds back into the expression for $A_r$ yields the desired result:
\begin{align*}
    A_r &\le -\frac{\eta\eta_r T}{4} + \frac{1}{8}\eta\eta_r T + \frac{1}{16}\eta\eta_r T \\
    &= \left( -\frac{4}{16} + \frac{2}{16} + \frac{1}{16} \right) \eta\eta_r T \\
    &= -\frac{1}{16}\eta\eta_r T.
\end{align*}
This completes the proof.
\end{proof}

\subsubsection{General Convergence Rate for Decaying Step-Sizes}

Applying the bound on $A_r$ from Lemma~\ref{lem:descent_conditions} to the main inequality \eqref{eq:telescoped-raw} yields:
\[
\E[f(x_R) - f(x_0)] \le \sum_{r=0}^{R-1} \left( -\frac{\eta\eta_r T}{16} \E\|\nabla f(w_r)\|^2 + E_r\,\E[\bar E_r] + V_r\,\sigma^{2} + H_r\,\nu^{2} \right).
\]
Rearranging the terms to isolate the sum of squared gradients and assuming the function is bounded below by $f_\star := \inf_x f(x)$, we obtain a bound on the weighted sum:

\begin{equation}
\label{eq:weighted_gradient_bound}
\sum_{r=0}^{R-1} \eta_r T\,\E\|\nabla f(w_r)\|^{2} \le \frac{16}{\eta}(f(x_0)-f_\star) + \frac{16}{\eta}\sum_{r=0}^{R-1}\E\!\bigl[E_r\,\bar E_r + V_r\,\sigma^{2} + H_r\,\nu^{2}\bigr].
\end{equation}


The main challenge is to bound the term involving the residual energy, $\sum_r E_r \E[\bar{E}_r]$. We first decouple the coefficient using the supremum $\mathcal{E}_{\max} := \sup_r E_r$. Then, we apply the bound on the sum of residuals from Lemma~\ref{lem:residual-A}, which is derived by unrolling the per-round recursion for $\bar{E}_r$:
\begin{align*}
\frac{16}{\eta} \sum_{r=0}^{R-1} E_r \E[\bar{E}_r] &\le \frac{16\mathcal{E}_{\max}}{\eta} \sum_{r=0}^{R-1} \E[\bar{E}_r] \\
&\le \frac{16\mathcal{E}_{\max}}{\eta(1-\rho_{\max})} \bar{E}_0 + \frac{16\mathcal{E}_{\max}(1-1/\delta)}{\eta(1-\rho_{\max})} \sum_{r=0}^{R-1} \E[B_r^{(\nabla)} + B_r^{(\nu,\sigma)}] \\
&\overset{(a)}{\le} \frac{16\mathcal{E}_{\max}(1-1/\delta)}{\eta(1-\rho_{\max})} \sum_{r=0}^{R-1} \E[B_r^{(\nabla)} + B_r^{(\nu,\sigma)}]
\end{align*}
where \textnormal{(a)} uses $\bar E_0=0$ since $e_0^{(k)}\equiv0$.

Substituting the residual recursion (Lemma~\ref{lem:residual-A}) into \eqref{eq:weighted_gradient_bound} introduces forcing terms.
Since \(B_r^{(\nabla)}\) is proportional to \(\|\nabla f(w_r)\|^2\),
it produces a contribution of the form
\(\Theta\sum_{r=0}^{R-1}\eta_r T\,\E\|\nabla f(w_r)\|^2\),
where we \emph{define}
\[ 
\begin{aligned}\label{eq:def-Theta}
\Theta
&:= \frac{16}{\eta}\times \frac{\mathcal{E}_{\max}}{1-\rho_{\max}}
\left(1-\frac{1}{\delta}\right)\,\beta^{2}\,
\sup_{0\le r<R}\frac{8\,\eta_r^{2}T^{2}+288\,L^{2}\eta_r^{4}T^{4}}{\eta_r T}\\
&\;=\;\frac{16}{\eta}\times \frac{\mathcal{E}_{\max}}{1-\rho_{\max}}
\left(1-\frac{1}{\delta}\right)\,\beta^{2}\,
\sup_{r}\bigl(8\,\eta_r T+288\,L^{2}\eta_r^{3}T^{3}\bigr).
\end{aligned}
\]
With this notation, the weighted inequality becomes
\begin{align}\label{eq:absorb-ineq} \notag
\bigl(1-\Theta\bigr)&\sum_{r=0}^{R-1}\eta_r T\,\E\|\nabla f(w_r)\|^2\\
&\le
\frac{16}{\eta}\Bigl[(f(x_0)-f_\star)
+\frac{\mathcal{E}_{\max}}{1-\rho_{\max}}\left(1-\frac{1}{\delta}\right)\,\sum_{r=0}^{R-1}\E B_r^{(\nu,\sigma)}
+\sum_{r=0}^{R-1}V_r\,\sigma^2+\sum_{r=0}^{R-1}H_r\,\nu^2\Bigr].
\end{align}
We now assume a small-steps/compression regime that ensures \(\Theta\le \tfrac12\). We can therefore move this term to the left-hand side and absorb it, which tightens the overall bound at the cost of multiplying the remaining terms on the right-hand side by a factor of 2.

To obtain a bound on the conventional average, compare it to the
weighted sum with weights \(q_r:=\eta_r T>0\). Let
\(q_{\min}:=\min_{0\le r<R} q_r\), \(S_R:=\sum_{r=0}^{R-1} q_r\), and
\(\displaystyle \phi_R:=\frac{S_R}{R\,q_{\min}}\ge 1\).
Then, for nonnegative terms,
\begin{align}
\frac{1}{R}\sum_{r=0}^{R-1}\E\|\nabla f(w_r)\|^{2}
&\le \frac{S_R}{R\,w_{\min}}\times \frac{1}{S_R}\sum_{r=0}^{R-1} w_r\,\E\|\nabla f(w_r)\|^{2}\\
&= \phi_R\times\frac{1}{S_R}\sum_{r=0}^{R-1}\eta_r T\,\E\|\nabla f(w_r)\|^{2}.
\end{align}
Combining this with the absorbed bound and noting \(x_0=w_0\) since \(e_0^{(k)}\equiv 0\), we obtain
\[
\begin{aligned}
\frac{1}{R}\sum_{r=0}^{R-1}\E\|\nabla f(w_r)\|^{2}
&\le\ \phi_R\,\frac{32}{\eta\,S_R}\Bigg[
(f(w_0)-f_\star)
+\sum_{r=0}^{R-1}V_r\,\sigma^{2}
+\sum_{r=0}^{R-1}H_r\,\nu^{2} \\[-2pt]
&\hspace{26mm}
+\ \frac{\mathcal E_{\max}}{1-\rho_{\max}}\Bigl(1-\tfrac1\delta\Bigr)
\sum_{r=0}^{R-1}\!\Bigl(4\,\eta_r^{2}T+96\,L^{2}\eta_r^{4}T^{3}\Bigr)\,\sigma^{2} \\[-2pt]
&\hspace{26mm}
+\ \frac{\mathcal E_{\max}}{1-\rho_{\max}}\Bigl(1-\tfrac1\delta\Bigr)
\sum_{r=0}^{R-1}\!\Bigl(8\,\eta_r^{2}T^{2}+1344\,L^{2}\eta_r^{4}T^{4}\Bigr)\,\nu^{2}
\Bigg].
\end{aligned}
\]

\subsubsection{Convergence Rate for a Constant Step-Size}

In the simpler setting where the inner learning rate is constant, $\eta_r \equiv \eta_0$ for all $r$, the sum of weights is $S_R = R\eta_0 T$. Under the descent and absorption conditions, we obtain
\begin{align*}
\frac{1}{R}\sum_{r=0}^{R-1}\E\|\nabla f(w_r)\|^{2}
\ \le\
&\frac{32\,(f(w_0)-f^\star)}{\eta\,\eta_0T\,R}
\;+\; \Bigl(1-\tfrac{1}{\delta}\Bigr)\Bigl[C_\sigma\,\eta_0^{2}L^{2}T\,\sigma^{2} + C_\nu\,\eta_0^{2}L^{2}T^{2}\,\nu^{2}\Bigr] \\[-2pt]
&\;+\; \frac{128\,L\,\eta_0}{K}\,\sigma^{2}
\end{align*}
with
\[
\begin{aligned}
C_\sigma&=\frac{32}{\eta}\Bigl[6\,\eta\,\eta_0^{2}L^{2}T+96\,L^{3}\eta\,\eta_0^{3}T^{2}\Bigr]
+\frac{32}{\eta}\times\frac{\mathcal E_{\max}}{1-\rho_{\max}}\Bigl[4\,\eta_0+96\,L^{2}\eta_0^{3}T^{2}\Bigr],\\
C_\nu&=\frac{32}{\eta}\Bigl[84\,\eta\,\eta_0^{2}L^{2}T^{2}+1344\,L^{3}\eta\,\eta_0^{3}T^{3}\Bigr]
+\frac{32}{\eta}\times\frac{\mathcal E_{\max}}{1-\rho_{\max}}\Bigl[8\,\eta_0T+1344\,L^{2}\eta_0^{3}T^{3}\Bigr].
\end{aligned}
\]
which proves the theorem.
\end{proof}

\subsection{Lemmas}

\begin{lemma}[Local-model drift under SA-PEF]      \label{lem:local-drift}
Fix a communication round \(r\) and an inner-loop horizon \(T\ge 1\).
Assume the step-size condition \(0 < \eta_r \le \frac{1}{8LT}\) and Assumptions \ref{A1}--\ref{A3}.

Then, for every local step \(t\in\{0,\dots,T-1\}\),
\begin{align}
\frac{1}{K} \sum_{k=1}^{K}
\mathbb{E}_r \left[
  \|w_{r+\frac12,t}^{(k)} - w_r\|^2
\right]
&\le
12\,\eta_r^{2}T\,\sigma^{2}
+168\,\eta_r^{2}T^{2}\,\nu^{2}
+36\,\eta_r^{2}T^{2}\,\beta^{2}\,\|\nabla f(w_r)\|^{2}\notag\\
&\quad
+3\,\alpha_r^{2}\times \frac{1}{K}\sum_{k=1}^{K}\E_r\|e_r^{(k)}\|^{2}.
\end{align}

\end{lemma}

\begin{proof}
All expectations are conditional on the randomness inside round \(r\).

Set \(u_{k,t} := w_{r+\frac12,t}^{(k)} - w_r\) (\(u_{k,0} = -\alpha_r e_r^{(k)}\)).
Client \(k\) updates by
\[
u_{k,t+1} = u_{k,t} - \eta_r
          \left(\nabla f_k(w_{r+\frac12,t}^{(k)}) + \xi_{k,t}\right),
\qquad
\mathbb{E}_r[\xi_{k,t}] = 0,\;
\mathbb{E}_r\|\xi_{k,t}\|^{2} \le \sigma^{2}.
\]

Using \(\|a - b\|^{2} = \|a\|^{2} - 2\langle a, b\rangle + \|b\|^{2}\), we get
\[
\mathbb{E}_r\|u_{k,t+1}\|^{2}
= \mathbb{E}_r\|u_{k,t}\|^{2}
   - 2\eta_r\,\mathbb{E}_r\langle u_{k,t}, \nabla f_k(w_{r+\frac12,t}^{(k)})\rangle
   + \eta_r^{2}\mathbb{E}_r\|\nabla f_k(w_{r+\frac12,t}^{(k)})\|^{2}
   + \eta_r^{2}\sigma^{2}.
\]

Write \(\nabla f_k(w_{r+\frac12,t}^{(k)}) = \nabla f_k(w_r) + \Delta_{k,t}\)
with \(\|\Delta_{k,t}\| \le L\|u_{k,t}\|\) (\(L\)-smoothness), so
\[
-2\eta_r\langle u_{k,t}, \Delta_{k,t} \rangle
    \le 2\eta_r L\|u_{k,t}\|^{2},
\]
and using Young’s inequality \( 2ab \le \gamma a^2 + \frac{1}{\gamma}b^2 \) with $\gamma = \frac{1}{4T}$,
\[
-2\eta_r\langle u_{k,t}, \nabla f_k(w_r) \rangle
\le \tfrac{1}{4T}\|u_{k,t}\|^{2} + 4\eta_r^{2}T\|\nabla f_k(w_r)\|^{2}.
\]

Since \(\eta_r \le 1/(8LT)\), then \(2\eta_r L \le 1/(4T)\), hence
\[
-2\eta_r\,\mathbb{E}_r\langle u_{k,t}, \nabla f_k(w_{r+\frac12,t}^{(k)})\rangle
\le \tfrac{1}{2T}\,\mathbb{E}_r\|u_{k,t}\|^{2}
        + 4\eta_r^{2}T\|\nabla f_k(w_r)\|^{2}.
\]

Also,
\[
\mathbb{E}_r\|\nabla f_k(w_{r+\frac12,t}^{(k)})\|^{2}
\le 2\|\nabla f_k(w_r)\|^{2} + 2L^{2}\mathbb{E}_r\|u_{k,t}\|^{2}.
\]

Combining gives:
\[
\mathbb{E}_r\|u_{k,t+1}\|^{2}
\le \left(1 + \tfrac{1}{2T} + 2\eta_r^{2}L^{2}\right)\mathbb{E}_r\|u_{k,t}\|^{2}
+ (2 + 4T)\eta_r^{2}\|\nabla f_k(w_r)\|^{2}
+ \eta_r^{2}\sigma^{2}.
\]

Let \(A := 1 + \tfrac{1}{2T} + 2\eta_r^{2}L^{2} \le 1 + \tfrac{1}{T} \le e^{1/T}\).
Define \(\bar D_t := \frac{1}{K} \sum_k \mathbb{E}_r\|u_{k,t}\|^{2}\).
Using Assumption~\ref{A3} we get
\[
\bar D_{t+1}
\le A\,\bar D_t
+ (2 + 4T)\eta_r^{2}
  \left(\beta^{2}\|\nabla f(w_r)\|^{2} + \nu^{2} \right)
+ \eta_r^{2}\sigma^{2} \;=:\; A\bar D_t + B.
\]

Since \(A \le e^{1/T}\), for \(t\le T\) we have \(A^{t} \le e \le 3\) and
\[
\frac{A^{t}-1}{A-1} \le \frac{e-1}{A-1} \le 2eT \le 6T,
\]
because \(A-1 \ge \tfrac{1}{2T}\).
Therefore,
\[
\bar D_t
\le A^{t}\bar D_0 + \frac{A^{t}-1}{A-1}\,B
\le 3\bar D_0
+ 6T \eta_r^{2} \sigma^{2}
+ 6T (2 + 4T) \eta_r^{2}
  \left( \beta^{2} \|\nabla f(w_r)\|^{2} + \nu^{2} \right).
\]

Finally, with \(\bar D_0 = \alpha_r^{2} \tfrac{1}{K} \sum_k \mathbb{E}_r \|e_r^{(k)}\|^{2}\), we loosen constants to the displayed form in the lemma:
\[
\bar D_t
\le
12\eta_r^{2}T
   \left(\sigma^{2} + 6T\nu^{2} + \beta^{2}\|\nabla f(w_r)\|^{2}\right)
+ 24\eta_r^{2}T^{2}
   \left(\beta^{2}\|\nabla f(w_r)\|^{2} + 4\nu^{2}\right)
+ 3\alpha_r^{2}\frac{1}{K} \sum_k \mathbb{E}_r \|e_r^{(k)}\|^{2}.
\]

Combine the $T$ and $T^2$ terms using $12T+24T^2\le 36T^2$ and compute $12T\times 6T+24T^2\times 4=168T^2$, we get
\begin{align}
\bar D_t
&\le
12\,\eta_r^{2}T\,\sigma^{2}
+168\,\eta_r^{2}T^{2}\,\nu^{2}
+36\,\eta_r^{2}T^{2}\,\beta^{2}\,\|\nabla f(w_r)\|^{2}
+3\,\alpha_r^{2}\times \frac{1}{K}\sum_{k=1}^{K}\E_r\|e_r^{(k)}\|^{2}.  \notag
\end{align}

\end{proof}



\begin{lemma}[Second moment of the shifted average update]
\label{lem:step-ahead-moment-fixed}
Let $K$ be the number of clients, $T\ge 1$ the number of local steps,
$\eta_r\in(0,\frac{1}{8LT}]$ the local stepsize in round $r$, and
$\alpha_r\in[0,1]$ the step-ahead parameter.
Define
\[
\bar e_r = \frac1K\sum_{k=1}^K e_r^{(k)},\qquad
g_r^{(k)} = \eta_r\sum_{t=0}^{T-1}
           \bigl(\nabla f_k(w_{r+\frac12,t}^{(k)})+\xi_{k,t}\bigr),\qquad
\bar g_r = \frac1K\sum_{k=1}^K g_r^{(k)}.
\]
Under Assumptions \ref{A1}--\ref{A3},
\begin{align}
\E_r\!\bigl[\|\alpha_r\bar e_r-\bar g_r\|^{2}\bigr]
&\le 
2\alpha_r^{2}\bar E_r + \;8\eta_r^{2}T^{2}\|\nabla f(w_r)\|^{2}\bigl(1+36L^{2}\eta_r^{2}T^{2}\beta^{2}\bigr) \notag\\
&\quad
\;+\;\frac{4\eta_r^{2}T\sigma^{2}}{K}+\;96L^{2}\eta_r^{4}T^{3}\bigl(\sigma^{2}+14T\nu^{2}\bigr)
\;+\;24\alpha_r^{2}\eta_r^{2}L^{2}T^{2}\bar E_r.
\end{align}

Here $\displaystyle
      \bar E_r = \frac1K\sum_{k=1}^K \|e_r^{(k)}\|^{2}$ is the average
residual energy at the beginning of round $r$, and $w_r$ is the global
model before local computation.
\end{lemma}

\begin{proof}
Using Young’s inequality, for any $a,b\in\R^d$, $\|a-b\|^2\le2\|a\|^2+2\|b\|^2$.  
With $a=\alpha_r\bar e_r$ (deterministic given~$\mathcal F_r$) 
and $b=\bar g_r$,
\begin{equation}
\|\alpha_r\bar e_r-\bar g_r\|^{2}
\le
2\alpha_r^{2}\|\bar e_r\|^{2}+2\|\bar g_r\|^{2}.
\label{eq:young-split}
\end{equation}
Since $\|\bar e_r\|^{2}\le\bar E_r$ (Jensen), 
the first term in \eqref{eq:young-split} gives $2\alpha_r^{2}\bar E_r$ after expectation.

Bounding the second moment of $\bar g_r$ and the gradient–noise cross term.
Write $\bar g_r=\eta_r\sum_{t=0}^{T-1}(\bar\nabla_t+\bar\xi_t)$ with
$\bar\nabla_t=\tfrac1K\sum_k\nabla f_k(w_{r+\frac12,t}^{(k)})$ and
$\bar\xi_t=\tfrac1K\sum_k\xi_{k,t}$.
By Young’s inequality,
\[
\E_r\|\bar g_r\|^{2}
\;\le\;
2\eta_r^{2}\,\E_r\Big\|\sum_{t=0}^{T-1}\bar\nabla_t\Big\|^{2}
\;+\;2\eta_r^{2}\,\E_r\Big\|\sum_{t=0}^{T-1}\bar\xi_t\Big\|^{2}.
\]
Using independence across $(k,t)$ and Assumption \ref{A2},
$\E_r\|\sum_t\bar\xi_t\|^{2}=\sum_t\E_r\|\bar\xi_t\|^{2}\le T\sigma^2/K$.

Bounding the summed local gradients.
Decompose
$\nabla f_k(w_{r+\frac12,t}^{(k)})=\nabla f_k(w_r)+\Delta_{k,t}$ with
$\|\Delta_{k,t}\|\le L\|w_{r+\frac12,t}^{(k)}-w_r\|$ (Assumption \ref{A1}). Then
\[
\frac{1}{K^{2}}
  \E_r\Bigl\|\sum_{k,t}\nabla f_k(w_{r+\frac12,t}^{(k)})\Bigr\|^{2}
\le
2T^{2}\|\nabla f(w_r)\|^{2}
+\frac{2L^{2}T}{K}
  \sum_{k,t}\E_r\|w_{r+\frac12,t}^{(k)}-w_r\|^{2}.
\]

Insert the local-model drift (Lemma~\ref{lem:local-drift}) and sum over $t$. 
For each $t\in\{0,\dots,T-1\}$, Lemma~\ref{lem:local-drift} yields
\[
\frac1K\sum_k\E_r\|w_{r+\frac12,t}^{(k)}-w_r\|^{2}
\le
12\,\eta_r^{2}T\,\sigma^{2}
+168\,\eta_r^{2}T^{2}\,\nu^{2}
+36\,\eta_r^{2}T^{2}\,\beta^{2}\,\|\nabla f(w_r)\|^{2} 
+3\alpha_r^2 \bar E_r.
\]
Summing over $t=0,\dots,T-1$ gives
\[
\frac1K\sum_{k,t}\E_r\|w_{r+\frac12,t}^{(k)}-w_r\|^{2}
\le
12\,\eta_r^{2}T^{2}\,\sigma^{2}
+168\,\eta_r^{2}T^{3}\,\nu^{2}
+36\,\eta_r^{2}T^{3}\,\beta^{2}\,\|\nabla f(w_r)\|^{2}
+3T\alpha_r^2 \bar E_r.
\]

Assemble,
\[
\begin{aligned}
\E_r\|\bar g_r\|^{2}
&\le
2\eta_r^{2}\Big[
2T^{2}\|\nabla f(w_r)\|^{2}
+\frac{2L^{2}T}{K}\!\sum_{k,t}\E_r\|w_{r+\frac12,t}^{(k)}-w_r\|^{2}
\Big]
+2\eta_r^{2}\times \frac{T\sigma^{2}}{K}\\
&\le
\;4\eta_r^{2}T^{2}\|\nabla f(w_r)\|^{2}\bigl(1+36L^{2}\eta_r^{2}T^{2}\beta^{2}\bigr)
\;+\;\frac{2\eta_r^{2}T\sigma^{2}}{K}+\;48L^{2}\eta_r^{4}T^{3}\bigl(\sigma^{2}+14T\nu^{2}\bigr)\\
&\quad
\;+\;12\alpha_r^{2}\eta_r^{2}L^{2}T^{2}\bar E_r.
\end{aligned}
\]
Finally, apply the factor \(2\) from \eqref{eq:young-split} and add
\(2\alpha_r^2\bar E_r\). This is precisely the assertion of the lemma.
\end{proof}





\begin{lemma}[Residual recursion under a $\delta$–contractive compressor]\label{lem:residual-A}
Fix a round $r$ and define $s_r:=\eta_r L T\le \tfrac18$. Let Assumption \ref{A1}--\ref{A3} hold and the compressor satisfies
Definition~\ref{def:compression-operator} with parameter $\delta\ge1$.
Let
\[
\bar E_r:=\frac1K\sum_{k=1}^K\|e_r^{(k)}\|^2,\qquad
g_r^{(k)}:=\eta_r\sum_{t=0}^{T-1}\bigl(\nabla f_k(w_{r+\frac12,t}^{(k)})+\xi_{k,t}\bigr).
\]
Then, with all expectations conditional on the randomness inside round $r$,
the averaged residual energy obeys
\begin{equation}\label{eq:resid-core}
\bar E_{r+1}\;\le\;\rho_r\,\bar E_r
\;+\;\Bigl(1-\tfrac1\delta\Bigr)\Bigl[B_r^{(\nabla)}+B_r^{(\nu,\sigma)}\Bigr],
\end{equation}
where
\[
\rho_r \;=\; \Bigl(1-\tfrac1\delta\Bigr)\Bigl(2(1-\alpha_r)^2+24\,\alpha_r^2 s_r^2\Bigr)
\;=\;\Bigl(1-\tfrac1\delta\Bigr)\Bigl(2-4\alpha_r+(2+24s_r^2)\alpha_r^2\Bigr),
\]
and
\[
\begin{aligned}
B_r^{(\nabla)}
&:=\ 8\,\eta_r^{2}T^{2}\,\beta^{2}\,\|\nabla f(w_r)\|^{2}
\;+\;288\,L^{2}\eta_r^{4}T^{4}\,\beta^{2}\,\|\nabla f(w_r)\|^{2},
\\[2pt]
B_r^{(\nu,\sigma)}
&:=\ 8\,\eta_r^{2}T^{2}\,\nu^{2}
\;+\;4\,\eta_r^{2}T\,\sigma^{2}
\;+\;96\,L^{2}\eta_r^{4}T^{3}\,\sigma^{2}
\;+\;1344\,L^{2}\eta_r^{4}T^{4}\nu^{2}.
\end{aligned}
\]
\end{lemma}

\begin{proof}
Write $u_{r+1}^{(k)}=(1-\alpha_r)e_r^{(k)}+g_r^{(k)}$ and
$e_{r+1}^{(k)}=u_{r+1}^{(k)}-C(u_{r+1}^{(k)})$.
By Definition~\ref{def:compression-operator},
\[
\E_r\|e_{r+1}^{(k)}\|^{2}\ \le\ \Bigl(1-\tfrac1\delta\Bigr)\,\E_r\|u_{r+1}^{(k)}\|^{2}.
\]
Average over $k$ and use $\|a+b\|^{2}\le 2\|a\|^{2}+2\|b\|^{2}$:
\[
\frac1K\sum_{k}\E_r\|u_{r+1}^{(k)}\|^{2}
\ \le\ 2(1-\alpha_r)^{2}\,\bar E_r\ +\ 2\times\frac1K\sum_{k}\E_r\|g_r^{(k)}\|^{2}.
\]
Next,
\[
\E_r\|g_r^{(k)}\|^{2}
\ \le\ 2\eta_r^{2}\,\E_r\Big\|\sum_{t}\nabla f_k(w_{r+\frac12,t}^{(k)})\Big\|^{2}
\ +\ 2\eta_r^{2}\,T\sigma^{2},
\]
and by $L$-smoothness,
\(
\nabla f_k(w_{r+\frac12,t}^{(k)})=\nabla f_k(w_r)+\Delta_{k,t}
\)
with
\(
\|\Delta_{k,t}\|\le L\|w_{r+\frac12,t}^{(k)}-w_r\|.
\)
Thus
\[
\E_r\Big\|\sum_{t}\nabla f_k(w_{r+\frac12,t}^{(k)})\Big\|^{2}
\ \le\ 2T^{2}\|\nabla f_k(w_r)\|^{2}\ +\ 2L^{2}T\sum_{t}\E_r\|w_{r+\frac12,t}^{(k)}-w_r\|^{2}.
\]
Averaging over $k$ and using the dissimilarity condition yields
\[
\frac1K\sum_{k}\E_r\|g_r^{(k)}\|^{2}
\ \le\ 4\eta_r^{2}T^{2}\bigl(\beta^{2}\|\nabla f(w_r)\|^{2}+\nu^{2}\bigr)
\ +\ 4\eta_r^{2}L^{2}T\,S
\ +\ 2\eta_r^{2}T\sigma^{2},
\]
where
\(
S:=\frac1K\sum_{k}\sum_{t=0}^{T-1}\E_r\|w_{r+\frac12,t}^{(k)}-w_r\|^{2}.
\)
By Lemma~\ref{lem:local-drift} (summed over \(t\)),
\[
S \le
12\,\eta_r^{2}T^{2}\sigma^{2}
+168\,\eta_r^{2}T^{3}\nu^{2}
+36\,\eta_r^{2}T^{3}\beta^{2}\|\nabla f(w_r)\|^{2}
+3T\alpha_r^2 \bar E_r.
\]
Substituting and noting $L^{2}\eta_r^{2}T^{2}=s_r^{2}$ gives
\[
\frac1K\sum_{k}\E_r\|g_r^{(k)}\|^{2}
\ \le\ \tfrac12\bigl(B_r^{(\nabla)}+B_r^{(\nu,\sigma)}\bigr)
\ +\ 12\,\alpha_r^{2}\,\eta_r^{2}L^{2}T^{2}\,\bar E_r .
\]
Insert this into the previous split to obtain
\[
\frac1K\sum_{k}\E_r\|u_{r+1}^{(k)}\|^{2}
\ \le\ \Bigl(2(1-\alpha_r)^{2}+24\,\alpha_r^{2}\eta_r^{2}L^{2}T^{2}\Bigr)\bar E_r
\ +\ B_r^{(\nabla)}+B_r^{(\nu,\sigma)}.
\]
Finally multiply by $(1-\tfrac1\delta)$ to conclude \eqref{eq:resid-core}.
\end{proof}

\begin{proposition}[Residual contraction vs.\ EF]\label{prop:sapef-rho}
Under the conditions of Lemma~\ref{lem:residual-A}, let
$\rho_{\mathrm{EF}}:=2\bigl(1-\tfrac1\delta\bigr)$ (the $\alpha_r=0$ baseline).
Then
\[
\rho_r
=\Bigl(1-\tfrac1\delta\Bigr)\Bigl(2-4\alpha_r+(2+24s_r^2)\alpha_r^2\Bigr),
\qquad s_r=\eta_rLT\le\tfrac18,
\]
and:
\begin{enumerate}
\item \emph{Strict improvement region.} For any
$\alpha_r\in\bigl(0,\tfrac{1}{\,1+12s_r^2\,}\bigr)$,
\[
\rho_r-\rho_{\mathrm{EF}}
=\Bigl(1-\tfrac1\delta\Bigr)\Bigl[-4\alpha_r+(2+24s_r^2)\alpha_r^2\Bigr]\;<\;0,
\]
and at the boundary $\alpha_r=\tfrac{2}{1+12s_r^2}$ the difference equals $0$.
\item \emph{Optimal step-ahead.}
The minimiser over $\alpha_r\in[0,1]$ is
\(
\alpha_r^\star=\frac{1}{1+12s_r^2}\in(0.84,1]
\)
and the minimum value is
\[
\rho_{\min}
=\Bigl(1-\tfrac1\delta\Bigr)\, \Bigl[2- \tfrac{2}{1+12s_r^2}\Bigr]
=\rho_{\mathrm{EF}}\Bigl(1-\tfrac{1}{1+12s_r^2}\Bigr).
\]
In particular, as $s_r\to 0$, $\rho_{\min}\to 0$ and
$\rho_{\min}/\rho_{\mathrm{EF}}\to 0$.
\end{enumerate}
\end{proposition}

\begin{proof}
Direct algebra from the quadratic form of $\rho_r$.
For (i), the sign of
$-4\alpha_r+(2+24s_r^2)\alpha_r^2
= \alpha_r\big((2+24s_r^2)\alpha_r-4\big)$
is negative when $\alpha_r<4/(2+24s_r^2)=1/(1+12s_r^2)$, and zero at equality.
For (ii), minimize
$q(\alpha)=2-4\alpha+(2+24s_r^2)\alpha^2$ to get
$\alpha_r^\star=2/(2+24s_r^2)=1/(1+12s_r^2)$ and
$q(\alpha_r^\star)=2-2/(1+12s_r^2)$; multiply by $(1-\tfrac1\delta)$.
\end{proof}

\subsection{Partial participation: analysis and rates}

At round $r$, let $\mathcal M_r\subseteq[K]$ be sampled uniformly without replacement,
$|\mathcal M_r|=m$, and denote $p:=m/K\in(0,1]$. Write $I_r^{(k)}=\1\{k\in\mathcal M_r\}$.
Active clients run the same inner loop as in full participation,
\[
g_r^{(k)}=\begin{cases}
\eta_r\sum_{t=0}^{T-1}\bigl(\nabla f_k(w_{r+\frac12,t}^{(k)})+\xi_{k,t}\bigr),& I_r^{(k)}=1,\\[2pt]
0,& I_r^{(k)}=0,
\end{cases}
\qquad
e_{r+1}^{(k)}=\begin{cases}
u_{r+1}^{(k)}-C(u_{r+1}^{(k)}),& I_r^{(k)}=1,\\[2pt]
e_r^{(k)},& I_r^{(k)}=0,
\end{cases}
\]
with $u_{r+1}^{(k)}=(1-\alpha_r)e_r^{(k)}+g_r^{(k)}$ for active clients. The server update is
\[
w_{r+1}=w_r-\eta\,\bar C_{r+1},\qquad
\bar C_{r+1}:=\frac{1}{m}\sum_{k\in\mathcal M_r}C\bigl(u_{r+1}^{(k)}\bigr).
\]
Expectations $\E_r[\cdot]$ are over local randomness and the draw of $\mathcal M_r$.
Assumptions \ref{A1}--\ref{A3} and Definition~\ref{def:compression-operator} (with $\delta\ge1$) hold.

\paragraph{Active/global averages; virtual iterate.}
Define
\[
\tilde e_r:=\frac{1}{K}\sum_{k=1}^{K}e_r^{(k)},\qquad
\bar e_r:=\frac{1}{m}\sum_{k\in\mathcal M_r}e_r^{(k)},\qquad
\bar g_r:=\frac{1}{m}\sum_{k\in\mathcal M_r}g_r^{(k)}.
\]
Note $\E_{\mathcal M_r}[\bar e_r\mid\{e_r^{(k)}\}]=\tilde e_r$. Let $x_r:=w_r-\eta\,\tilde e_r$.

\begin{lemma}[PP virtual-iterate identity]\label{lem:pp-virt}
With the definitions above,
\begin{equation}\label{eq:PP-virt}
x_{r+1}-x_r
=\eta\Bigl[\,p\bigl(\alpha_r\bar e_r-\bar g_r\bigr)\;-\;(1-p)\,\bar C_{r+1}\Bigr].
\end{equation}
\end{lemma}

\begin{lemma}[PP compression second moment]\label{lem:pp-Cbar}
Let $c_\delta:=2\bigl(2-\tfrac1\delta\bigr)$. Then, conditionally on $\mathcal M_r$,
\[
\E_r\|\bar C_{r+1}\|^{2}\;\le\;\frac{1}{m}\sum_{k\in\mathcal M_r}\E_r\|C(u_{r+1}^{(k)})\|^{2}
\;\le\;\frac{c_\delta}{m}\sum_{k\in\mathcal M_r}\E_r\|u_{r+1}^{(k)}\|^{2}.
\]
\end{lemma}

\begin{lemma}[One-round descent under PP]\label{lem:pp-descent}
Let $\Delta_r^{\mathrm{act}}:=\alpha_r\bar e_r-\bar g_r$. Under the same alignment and
stepsize coupling as in full participation,
\[
\text{\emph{(C1)}}\ \ \eta_r^{2}L^{2}\beta^{2}(6T^{2}+12T^{3})\le \tfrac{T}{8},
\qquad
\text{\emph{(C2)}}\ \ \eta\le \frac{1}{256\,\beta^{2}\,L\,\eta_rT},
\]
there exists a universal $c>0$ such that
\begin{equation}\label{eq:pp-per-round}
\E_r\bigl[f(x_{r+1})-f(x_r)\bigr]
\ \le\
-\,c\,p\,\eta\,\eta_rT\,\|\nabla f(w_r)\|^{2}
\ +\ \E_r\bigl[\widehat{\mathcal E}_r\,\tilde E_r\bigr]
\ +\ \widehat{\mathcal V}_r\,\sigma^{2}
\ +\ \widehat{\mathcal H}_r\,\nu^{2},
\end{equation}
where $\tilde E_r:=\frac{1}{K}\sum_{k=1}^{K}\|e_r^{(k)}\|^{2}$ and
$\widehat{\mathcal E}_r,\widehat{\mathcal V}_r,\widehat{\mathcal H}_r$
equal the full-participation coefficients scaled by $p$ and augmented by
$(1-p)$–terms originating from $\bar C_{r+1}$ via Lemma~\ref{lem:pp-Cbar}.
\end{lemma}

\begin{lemma}[Residual recursion under PP]\label{lem:pp-resid}
Let
\[
\rho_r^{\mathrm{PP}}:=(1-p)\;+\;p\Bigl(1-\tfrac1\delta\Bigr)\Bigl(2(1-\alpha_r)^2+24\,\alpha_r^2(\eta_rLT)^2\Bigr).
\]
Then
\begin{equation}\label{eq:pp-resid}
\E_r \tilde E_{r+1}
\ \le\ \rho_r^{\mathrm{PP}}\ \tilde E_r
\ +\ p\Bigl(1-\tfrac1\delta\Bigr)\Bigl[B_r^{(\nabla)}+B_r^{(\nu,\sigma)}\Bigr],
\end{equation}
with $B_r^{(\nabla)}$ and $B_r^{(\nu,\sigma)}$ as in full participation.
\end{lemma}

\paragraph{Telescoping, absorption, and final bounds (PP).}
Let $S_R:=\sum_{r=0}^{R-1}\eta_rT$ and $S_R^{\mathrm{PP}}:=\sum_{r=0}^{R-1}p\,\eta_rT=p\,S_R$.
Summing \eqref{eq:pp-per-round} over $r$ and using $\mathcal C_r^{\mathrm{PP}}\le -c\,p\,\eta\,\eta_rT$,
\[
\sum_{r=0}^{R-1}p\,\eta_rT\;\E\|\nabla f(w_r)\|^{2}
\ \le\ \frac{16}{\eta}\bigl(f(x_0)-\E f(x_R)\bigr)
+\frac{16}{\eta}\sum_{r=0}^{R-1}\E\bigl[\widehat{\mathcal E}_r\,\tilde E_r
+\widehat{\mathcal V}_r\,\sigma^{2}+\widehat{\mathcal H}_r\,\nu^{2}\bigr].
\]
If $f$ is bounded below by $f_\star$, then $f(x_R)\ge f_\star$. Summing \eqref{eq:pp-resid} and
assuming $\rho_{\max}^{\mathrm{PP}}:=\sup_r\rho_r^{\mathrm{PP}}<1$,
\[
\sum_{r=0}^{R-1}\E\tilde E_r
\ \le\ \frac{1}{1-\rho_{\max}^{\mathrm{PP}}}\,\tilde E_0
+\frac{p}{1-\rho_{\max}^{\mathrm{PP}}}\Bigl(1-\tfrac1\delta\Bigr)
\sum_{r=0}^{R-1}\E\bigl[B_r^{(\nabla)}+B_r^{(\nu,\sigma)}\bigr].
\]
Plugging this into the previous display yields a gradient-forcing term proportional to
\[
\frac{16}{\eta}\times\frac{\widehat{\mathcal E}_{\max}}{1-\rho_{\max}^{\mathrm{PP}}}\Bigl(1-\tfrac1\delta\Bigr)\times
\underbrace{p\sum_r B_r^{(\nabla)}}_{\le\ \beta^2 d_{\max}\,p\sum_r\eta_rT\,\E\|\nabla f(w_r)\|^2},
\]
where $d_{\max}:=\sup_r\bigl(8\,\eta_rT+288\,L^{2}\eta_r^{3}T^{3}\bigr)$. Defining
\[
\Theta_{\mathrm{PP}}
:= \frac{16}{\eta}\times\frac{\widehat{\mathcal E}_{\max}}{1-\rho_{\max}^{\mathrm{PP}}}\Bigl(1-\tfrac1\delta\Bigr)\beta^{2} d_{\max},
\]
we obtain the absorption inequality (the factor $p$ cancels on both sides):
\[
\bigl(1-\Theta_{\mathrm{PP}}\bigr)\sum_{r=0}^{R-1}p\,\eta_rT\,\E\|\nabla f(w_r)\|^{2}
\ \le\ \frac{16}{\eta}\Bigl[(f(x_0)-f_\star)+\sum_{r=0}^{R-1}\widehat{\mathcal V}_r\,\sigma^{2}+\widehat{\mathcal H}_r\,\nu^{2}\bigr]\Bigr].
\]
Assuming $\Theta_{\mathrm{PP}}\le\frac12$, we conclude
\[
\sum_{r=0}^{R-1}p\,\eta_rT\,\E\|\nabla f(w_r)\|^{2}
\ \le\ \frac{32}{\eta}\Bigl[(f(x_0)-f_\star)+\sum_{r=0}^{R-1}\widehat{\mathcal V}_r\,\sigma^{2}+\widehat{\mathcal H}_r\,\nu^{2}\Bigr].
\]
Dividing by $S_R^{\mathrm{PP}}=p\,\eta_0TR$ and with $\eta_r\equiv\eta_0$ yields the averaged bounds below.
\begin{align}
\frac{1}{R}\sum_{r=0}^{R-1}\E\|\nabla f(w_r)\|^{2}
~\le~&
~\frac{32}{\eta\,p\,\eta_0 T R}\,(f(x_0)-f_\star)
~+~ \frac{128\,L\,\eta_0}{\eta\,p\,m}\,\sigma^{2}
\nonumber\\
&~+~ \frac{32}{\eta\,p}\Bigl(1-\tfrac{1}{\delta}\Bigr)\Bigl[C_\sigma\,\eta_0^{2} L^{2} T\,\sigma^{2}+C_\nu\,\eta_0^{2} L^{2} T^{2}\,\nu^{2}\Bigr].
\label{eq:const-eta0-bound}
\end{align}

So, the optimization term scales as
$O\big((p\,\eta\,\eta_0TR)^{-1}\big)$ (a per-round slow-down by $1/p=K/m$).
The \emph{pure mini-batch} variance enjoys a $1/m$ reduction: its contribution scales as $\frac{128\,L\,\eta_0}{m}\,\sigma^{2}$,
while the residual-induced variance/heterogeneity floors scale as
$\Bigl(1-\tfrac1\delta\Bigr)\!\bigl[C_\sigma\,\eta_0^{2}L^{2}T\,\sigma^{2}
+ C_\nu\,\eta_0^{2}L^{2}T^{2}\nu^{2}\bigr]$.

\paragraph{Stalling vs.\ step-ahead.}
With $\alpha_r=0$ (no step-ahead), the multiplicative factor becomes
$\rho_r^{\mathrm{PP}}=(1-p)+2p(1-\tfrac1\delta)=1+p(1-\tfrac{2}{\delta})$, which can be
$\ge1$ under aggressive compression and small $p$, explaining the slowdown in cross-device regimes. For moderate $\alpha_r$ (e.g., $\alpha_r\approx\alpha_r^\star$ from the full-participation analysis), $\rho_r^{\mathrm{PP}}$ strictly decreases, improving the decay of $\tilde E_r$ each time a client participates and restoring faster early progress.

\paragraph{Constants and feasibility.}
The constants $C_\sigma$, $C_\nu$ and $\Theta$ in Theorem~\ref{main_theorem} collect the contributions of stochastic-gradient variance, data heterogeneity, and compression–induced residual drift. Inspecting their explicit formulas, we see that they depend on the algorithmic hyperparameters only through the effective local stepsize
\[
s_0 \;=\; \eta_0 L T,
\]
the compression bias factor $(1 - 1/\delta)$, and the residual–contraction term $1/(1-\rho_{\max})$. In particular, $\rho_{\max}$ itself is an increasing function of $s_0$ and $(1-1/\delta)$, so $C_\sigma$, $C_\nu$ and $\Theta$ are \emph{monotone nondecreasing} in $s_0$ and $(1-1/\delta)$: larger local work or more aggressive compression lead to larger constants and thus a higher residual-driven floor.

In the partial-participation extension (Remark~\ref{rem:pp}), the corresponding constant $\Theta_{\mathrm{PP}}$ inherits the same monotone dependence on $s_0$ and $(1-1/\delta)$ and, in addition, scales inversely with the participation rate $p = m/K$ (i.e., it increases as $p$ decreases). Thus, the feasibility conditions $\rho_{\max}^{\mathrm{PP}} < 1$ and $\Theta_{\mathrm{PP}} \le \tfrac12$ can be interpreted as requiring a standard “small” effective local stepsize $s_0$, moderate compression, and not-too-extreme partial participation. For the default hyperparameters used in our experiments (e.g., $T=5$, Top-$k$ compression, and $p \in \{0.1,0.2,1.0\}$), we numerically evaluate these constants and confirm that $\rho_{\max}^{\mathrm{PP}} < 1$ in all regimes. Moreover, in a mildly compressed setting (e.g., $\delta = 1.005$ with the same stepsizes), the corresponding $\Theta_{\mathrm{PP}}$ lies well below $\tfrac12$, illustrating that the condition $\Theta_{\mathrm{PP}} \le \tfrac12$ is a conservative sufficient condition rather than a tight practical tuning rule for the aggressively compressed Top-$k$ regimes we study.

% \section{Useful Inequalities}

% We collect a few inequalities used throughout the proofs.

% \begin{lemma}[Sum vs.\ sum of squares]\label{lem:sum-bound}
% For any vectors $a_1,\dots,a_n\in\mathbb{R}^d$,
% \[
% \Big\|\sum_{i=1}^n a_i\Big\|^2 \;\le\; n \sum_{i=1}^n \|a_i\|^2.
% \]
% \end{lemma}

% \begin{lemma}[Young’s inequality for products]\label{lem:young}
% For any $a,b\in\mathbb{R}^d$ and any $\gamma>0$,
% \[
% 2\langle a,b\rangle \;\le\; \gamma\,\|a\|^2 \;+\; \gamma^{-1}\,\|b\|^2.
% \]
% \end{lemma}

% \begin{lemma}[Two-term quadratic bound]\label{lem:two-term}
% For any $a,b\in\mathbb{R}^d$ and any $\alpha>0$,
% \[
% \|a+b\|^2 \;\le\; (1+\alpha)\,\|a\|^2 \;+\; (1+\alpha^{-1})\,\|b\|^2.
% \]
% \end{lemma}

\subsection{Comparison of EF-type compressed FL methods.}
For completeness, In Table~\ref{tab:rate-comparison}, we summarizes the main assumptions, mechanisms, and qualitative nonconvex behavior of several closely related algorithms: Fed-EF, SAEF, CSER, SCAFCOM, and SA-PEF. The goal is not to restate full theorems, but to highlight the regimes they target. Fed-EF and SA-PEF share the same lightweight, stateless EF architecture and standard FL assumptions (local steps, partial participation, biased contractive compressors), with SA-PEF improving the residual-contraction constant under compression. CSER and SAEF focus on centralized/local-SGD settings without partial participation, while SCAFCOM achieves stronger robustness to arbitrary heterogeneity by adding SCAFFOLD-style control variates and momentum, at the cost of increasing per-client state and communication.

\begin{table}[t]
\centering
\scriptsize
\caption{Comparison of compressed algorithms for FL. \textbf{SA-PEF} bridges the gap between Fed-EF (stable but slower under aggressive compression) and SAEF (faster warm-up but fragile in heterogeneous FL), achieving improved residual contraction without the extra state and complexity of control-variate methods such as SCAFCOM.}
\setlength{\tabcolsep}{3pt}
\begin{tabular}{p{1.2cm} p{2.1cm} p{2.6cm} p{7.22cm}}
\toprule
\textbf{Algorithm}
& \textbf{Assumptions \newline(beyond\newline $L$-smoothness)}
& \textbf{Mechanism \& State}
& \textbf{Convergence / Behaviour (Nonconvex)} \\
\midrule
Fed-EF\newline \citep{li2023analysis}
& Bounded variance; local SGD with PP
& Biased $\delta$-contractive; \newline \textbf{Stateless}\newline (one residual/client)
& Standard FL nonconvex rate with $1/p$ slowdown under partial participation; 
residual contraction factor $\rho_{\text{EF}}$ can be relatively weak under aggressive compression, leading to slower progress and earlier stalling. \\
[0.4em]
SAEF\newline \citep{xu2021step}
& Bounded gradient;\newline centralized, synchronous setting;\newline no PP analysis
& Full step-ahead ($\alpha=1$); \newline \textbf{Stateless}\newline (one residual/client)
& Analyzed in the classical distributed setting (no local steps, no client sampling); reduces EF’s gradient mismatch there. \newline
In our FL experiments with heterogeneous data, full step-ahead tends to produce larger gradient mismatch and late-stage plateaus compared to EF/SA-PEF (Sec.~4). \\
[0.4em]
CSER\newline \citep{xie2020cser}
& Bounded variance; local SGD (typically full participation);
& Error reset\newline (periodic dense communication of residuals); \newline \textbf{Stateless}
& Controls residual drift via periodic resets, yielding a nonconvex local-SGD rate with an $R$-independent floor. \newline
However, resets require sending full residuals, inducing \textbf{high peak bandwidth} at reset rounds and analyses usually assume full participation. \\
[0.4em]
SCAFCOM\newline \citep{huang2024stochastic}
& \textbf{Arbitrary heterogeneity}; bounded variance; local SGD with PP
& Control variates\newline + momentum; \newline \textbf{Stateful}\newline (extra state $c$,$c_i$ per client)
& Achieves a nonconvex FL guarantees under arbitrary non-IID data, with improved dependence on heterogeneity, \newline
at the cost of \textbf{higher system complexity} (maintaining control-variates state). \\
[0.4em]
\textbf{SA-PEF (ours)}
& Gradient dissimilarity $(\beta,\nu)$; local SGD with PP
& Partial step-ahead\newline ($0<\alpha<1$); \newline \textbf{Stateless}\newline (one residual/client)
& Operates in the same FL regime and under the same assumptions as Fed-EF, with the same leading-order nonconvex rate, \newline
but with a \textbf{strictly improved} residual contraction factor $\rho_{\max} < \rho_{\text{EF}}$ under biased compression, which lowers the error floor and balances warm-up speed with long-term stability. \\
\bottomrule
\end{tabular}
\label{tab:rate-comparison}
\end{table}

\begin{remark}[Relation to EF21]
EF21 and its extensions~\citep{richtarik2021ef21,fatkhullin2025ef21bells} obtain stronger guarantees (no error floor) under a different regime: synchronized data-parallel training with $T{=}1$, full-gradient (or gradient-difference) compression at a shared iterate, and no local steps. In this setting, EF21 is strictly preferable to classical EF. Our analysis targets a complementary regime, federated local-SGD with $T>1$ local steps, partial participation, and biased contractive compressors, where the compressed object is the accumulated local update and client drift plays a central role. Extending EF21-style arguments (or designing EF21-style step-ahead variants) to this local-SGD, partial-participation, biased-compression setting is non-trivial and remains an interesting direction for future work.
\end{remark}

\section{Implementation Details and Additional Experiments}
\subsection{Setup and Parameter tuning}
This appendix details datasets, federated partitioning, compressors, the hyperparameter search protocol, and fairness controls used across all methods.

\paragraph{Datasets, models, and preprocessing.}
{CIFAR-10 (ResNet-9)}, {CIFAR-100 (ResNet-18)}, and {Tiny-ImageNet} (64$\times$64; {ResNet-34}) are trained with cross-entropy loss.
Preprocessing follows standard practice: per-dataset mean/std normalization; CIFAR uses random crop (4-pixel padding) and horizontal flip; Tiny-ImageNet uses random resized crop to 64 and horizontal flip. Unless stated otherwise, batch size is 64, momentum is 0.9, weight decay is $5\!\times\!10^{-4}$ on CIFAR and $10^{-4}$ on Tiny-ImageNet.

\paragraph{Federated partitioning, participation, and local computation.}
We create $K{=}100$ clients and apply Dirichlet label partitioning with $\gamma\in\{0.1,1.0\}$ (smaller $\gamma$ $\Rightarrow$ stronger non-IID).
Each round samples $m=\lfloor pK\rfloor$ clients uniformly without replacement with $p\in\{0.1,0.5,1.0\}$.
Participating clients run $T$ local SGD steps at stepsize $\eta_r$; default $T{=}5$. We train for $R{=}200$ rounds.
Unless stated, server stepsize is $\eta{=}1.0$. We conducted all federated learning simulations using the \textsc{Flower} framework~\citep{beutel2020flower}.

\paragraph{Compressors and communication accounting.}
We use Top-$k$ sparsification with $k/d\in\{0.01,0.05,0.10\}$; each selected entry communicates its \emph{index} and \emph{value}. As a consequence, Top-$k$ satisfies Definition~\ref{def:compression-operator} with $\delta=d/k$; we record the standard bound below.

\begin{lemma}[Top-$k$ contraction; \citep{stich2018sparsified,beznosikov2023biased}]
\label{lem:topk}
Let $C=\operatorname{Top}_k:\mathbb{R}^d\!\to\!\mathbb{R}^d$ keep the $k$ largest
absolute-value coordinates of $x$ (ties broken arbitrarily), zeroing the rest.
Then for all $x\in\mathbb{R}^d$,
\[
\|x-C(x)\|_2^2 \;\le\; \Bigl(1-\frac{k}{d}\Bigr)\,\|x\|_2^2,
\qquad
\frac{k}{d}\,\|x\|_2^2 \;\le\; \|C(x)\|_2^2 \;=\; \langle C(x),x\rangle \;\le\; \|x\|_2^2.
\]
In particular, $C$ is $\delta$-contractive with $\delta=d/k$ in the sense of
Definition~\ref{def:compression-operator}. The constants are tight when all $|x_i|$
are equal.
\end{lemma}


Each selected entry transmits its \emph{index} and \emph{value}. \emph{Raw uplink bits} per participating client per round are
$k\big(\lceil\log_2 d\rceil + b_{\text{val}}\big)$ with $b_{\text{val}}{=}32$ for FP32 values. Unless stated, reported \emph{cumulative communication} aggregates \emph{uplink only} across participating clients (downlink is identical across compressed methods and omitted for fairness); FedAvg’s downlink/uplink are both dense FP32.

\paragraph{Hyperparameter search protocol.}
We adopt a small, method-agnostic grid tuned on a held-out validation split. Unless noted, we select hyperparameters by best \emph{validation top-1} at a \emph{fixed communication budget} (bits) within $R$ rounds; ties are broken by higher accuracy at earlier checkpoints.
We reuse the \emph{same} grid across participation rates ($q$).

\textbf{Search spaces (shared across methods).}
\label{search_grid}
\begin{itemize}[leftmargin=8pt]
\item \textbf{Client LR $\eta_r$:}
\begin{tabular}{@{}l l@{}}
CIFAR-10:      & $\{0.001, 0.05, 0.1, 0.2, 1.0, 10\}$ \\
CIFAR-100:     & $\{0.001, 0.05, 0.1, 1.0, 10\}$ \\
Tiny-ImageNet: & $\{0.001, 0.02, 0.05, 1.0\}$ \\
               & (cosine decay with 3-5 epoch warm-up, minimum $\eta_r = 0.005$)
\end{tabular}
\item \textbf{Server LR $\eta$:} $\{0.5,1.0\}$.
\item \textbf{Weight decay:} $\{5\!\times\!10^{-4},\,10^{-4}\}$.
\item \textbf{Local steps $T$:} $\{1,5,10\}$.
\item \textbf{Compressor level $k/d$:} $\{0.01,0.05,0.10\}$.
\end{itemize}
\textbf{SA-PEF-specific.} Constant-$\alpha$ ablations use $\alpha\in\{0.0,0.1, 0.2, 0.3, 0.4, 0.5, 0.6, 0.7, 0.8, 0.9, 1.0\}$; unless stated, $\alpha$ is fixed across rounds.

\paragraph{Fairness controls and evaluation.}
(i) The \emph{same} grid is used across methods; (ii) the best setting is selected at a matched bit budget; (iii) client sampling seeds are shared across methods; (iv) evaluation uses the server model in \texttt{eval} mode with identical preprocessing.
We report both \emph{accuracy vs.\ rounds} and \emph{accuracy vs.\ bits}; the latter is our primary metric under communication constraints. Experiments ran on NVIDIA A100/A5000/H200 GPUs; hardware does not affect communication accounting.


\subsection{Additional Experiments}
\paragraph{Multi-seed stability.} In Table~\ref{tab:error_bar}, we report final test accuracy as mean $\pm$ standard deviation over five independent runs with different random seeds. Across seeds, the relative ranking of methods is consistent, with SA-PEF retaining its advantage in high-compression, low-participation regimes.

\begin{table}[t]
\centering
\caption{Final test accuracy (mean $\pm$ std.\ over five independent runs) for CIFAR-10 and CIFAR-100 under two participation/heterogeneity regimes. The hyperparameters used in all algorithms are $R=200$, $T=5$, and $\eta_l=0.1$.}
\label{tab:error_bar}
%\setlength{\tabcolsep}{3pt}
\begin{tabular}{lllcccc}
\toprule
\multirow{3}{*}{\textbf{Dataset}} & \multirow{3}{*}{\textbf{Model}} & \multirow{3}{*}{\textbf{Algorithm}} & \multicolumn{4}{c}{\textbf{Final test accuracy (\%)}} \\
 & & & \multicolumn{2}{c}{$p = 0.1,\ \gamma = 0.5$} & \multicolumn{2}{c}{$p = 0.5,\ \gamma = 0.1$} \\
 & & & top-1 & top-10 & top-1 & top-10 \\
\midrule
\multirow{5}{*}{CIFAR-10} & \multirow{5}{*}{ResNet-9} 
  & FedAvg  & \multicolumn{2}{c}{$88.5 \pm 1.6$} & \multicolumn{2}{c}{$69.5 \pm 1.5$} \\
 & & EF      & $57.0 \pm 2.0$ & $72.7 \pm 2.3$ & $40.3 \pm 3.7$ & $47.2 \pm 4.1$ \\
 & & SAEF    & $74.2 \pm 3.9$ & $75.5 \pm 2.9$ & $39.5 \pm 4.6$ & $48.5 \pm 3.3$ \\
 & & CSER    & $68.6 \pm 2.8$ & $80.5 \pm 1.2$ & $49.5 \pm 4.0$ & $67.2 \pm 2.6$ \\
 & & SA-PEF  & $80.5 \pm 2.6$ & $82.7 \pm 1.6$ & $47.5 \pm 3.6$ & $68.5 \pm 1.9$ \\
\midrule
\multirow{5}{*}{CIFAR-100} & \multirow{5}{*}{ResNet-18} 
  & FedAvg  & \multicolumn{2}{c}{$62.5 \pm 1.6$} & \multicolumn{2}{c}{$61.9 \pm 0.6$} \\
 & & EF      & $40.4 \pm 2.4$ & $49.5 \pm 1.6$ & $41.5 \pm 1.6$ & $57.5 \pm 1.9$ \\
 & & SAEF    & $46.1 \pm 1.4$ & $50.5 \pm 2.0$ & $44.5 \pm 3.6$ & $57.5 \pm 3.9$ \\
 & & CSER    & $46.2 \pm 0.4$ & $51.5 \pm 1.9$ & $42.5 \pm 2.4$ & $60.7 \pm 1.0$ \\
 & & SA-PEF  & $49.6 \pm 1.8$ & $54.5 \pm 2.6$ & $48.5 \pm 2.0$ & $60.6 \pm 2.8$ \\
\bottomrule
\end{tabular}
\end{table}

Figures \ref{fig:cifar10_resnet9_100_100}--\ref{fig:Tiny-ImageNet_resnet34_100_100_5} report extra convergence curves for SA-PEF and the baselines on \textbf{CIFAR-10}, \textbf{CIFAR-100}, and \textbf{Tiny-ImageNet}. For each dataset we sweep (i) \emph{participation} \(q\in\{1.0,\,0.1\}\), (ii) \emph{compression budget} (Top-1\% and Top-10\% under full participation; Top-5\% and Top-1\% under \(q{=}0.1\)), and (iii) \emph{data heterogeneity} via Dirichlet partitions. The plots include both \emph{accuracy vs.\ rounds} and \emph{accuracy vs.\ communicated GB}. Under full participation, SA-PEF consistently reaches a given accuracy in fewer rounds than EF/CSER and tracks SAEF without late-stage plateaus. Under partial participation with aggressive compression (Top-5\%, Top-1\%), the gaps naturally narrow but the qualitative trend persists, illustrating that the main conclusions are robust across datasets, architectures, and federation settings.

\begin{figure*}[htbp]
  \centering
  % Row 1
  \begin{subfigure}{0.45\textwidth}
    \includegraphics[width=\linewidth]{Plots/efficiency_plt_100_participation_top_1_resnet9.pdf}
    \caption{$p {=} 1.0$, Top-1\%}\label{fig:cifar100_100_1_e}
  \end{subfigure}\hfill
  \begin{subfigure}{0.45\textwidth}
    \includegraphics[width=\linewidth]{Plots/efficiency_plt_100_participation_top_10_resnet9.pdf}
    \caption{$p {=} 1.0$, Top-10\%}\label{fig:cifar100_100_10_e}
  \end{subfigure}
  \vspace{4pt}
  % Row 2
  \begin{subfigure}{0.45\textwidth}
    \includegraphics[width=\linewidth]{Plots/plt_100_participation_top_1_resnet9.pdf}
    \caption{$p {=} 1.0$, Top-1\%}\label{fig:cifar100_100_1_c}
  \end{subfigure}\hfill
  \begin{subfigure}{0.45\textwidth}
    \includegraphics[width=\linewidth]{Plots/plt_100_participation_top_10_resnet9.pdf}
    \caption{$p {=} 1.0$, Top-10\%}\label{fig:cifar100_100_10_c}
  \end{subfigure}
  \caption{Test accuracy vs number of rounds (row 1) and communicated GB (row 2) on the CIFAR-10 dataset using ResNet-9 and $\gamma{=}0.1$.}% across participation rates with Top-10\%
  \label{fig:cifar10_resnet9_100_100}
\end{figure*}


\begin{figure*}[htbp]
  \centering
  % Row 1
  \begin{subfigure}{0.45\textwidth}
    \includegraphics[width=\linewidth]{Plots/efficiency_plt_10_participation_top_5_resnet18.pdf}
    \caption{$p {=} 0.1$, Top-5\%}\label{fig:cifar100_10_5_e}
  \end{subfigure}\hfill
  \begin{subfigure}{0.45\textwidth}
    \includegraphics[width=\linewidth]{Plots/efficiency_plt_50_participation_top_5_resnet18.pdf}
    \caption{$p {=} 0.5$, Top-5\%}\label{fig:cifar100_50_5_e}
  \end{subfigure}
  \vspace{4pt}
  % Row 2
  \begin{subfigure}{0.45\textwidth}
    \includegraphics[width=\linewidth]{Plots/plt_10_participation_top_5_resnet18.pdf}
    \caption{$p {=} 0.1$, Top-5\%}\label{fig:cifar100_10_5_c}
  \end{subfigure}\hfill
  \begin{subfigure}{0.45\textwidth}
    \includegraphics[width=\linewidth]{Plots/plt_50_participation_top_5_resnet18.pdf}
    \caption{$p {=} 0.5$, Top-5\%}\label{fig:cifar100_50_5_c}
  \end{subfigure}
  \caption{Test accuracy vs number of rounds (row 1) and communicated GB (row 2) on the CIFAR-100 dataset using ResNet-18 and $\gamma{=}0.1$.}% across participation rates with Top-10\%
  \label{fig:cifar10_resnet9_100_100_5}
\end{figure*}

\begin{figure*}[htbp]
  \centering
  % Row 1
  \begin{subfigure}{0.48\textwidth}
    \includegraphics[width=\linewidth]{Plots/efficiency_plt_10_participation_top_1_resnet9.pdf}
    \caption{$p {=} 0.1$, Top-1\%}\label{fig:cifar10_10_1_e}
  \end{subfigure}\hfill
  \begin{subfigure}{0.48\textwidth}
    \includegraphics[width=\linewidth]{Plots/efficiency_plt_10_participation_top_10_resnet9.pdf}
    \caption{$p {=} 0.1$, Top-10\%}\label{fig:cifar10_10_10_e}
  \end{subfigure}
  \vspace{4pt}
  % Row 2
  \begin{subfigure}{0.48\textwidth}
    \includegraphics[width=\linewidth]{Plots/plt_10_participation_top_1_resnet9.pdf}
    \caption{$p {=} 0.1$, Top-1\%}\label{fig:cifar10_10_1_c}
  \end{subfigure}\hfill
  \begin{subfigure}{0.48\textwidth}
    \includegraphics[width=\linewidth]{Plots/plt_10_participation_top_10_resnet9.pdf}
    \caption{$p {=} 0.1$, Top-10\%}\label{fig:cifar10_10_10_c}
  \end{subfigure}
  \caption{Test accuracy vs number of rounds (row 1) and communicated GB (row 2) on the CIFAR-10 dataset using ResNet-9 and $\gamma{=}0.1$.}% across participation rates with Top-10\%
  \label{fig:cifar10_resnet9}
\end{figure*}

\begin{figure*}[htbp]
  \centering
  \begin{subfigure}{0.45\textwidth}
    \includegraphics[width=\linewidth]{Plots/efficiency_plt_10_participation_top_10_resnet34_iid.pdf}
  \end{subfigure}\hfill
  \begin{subfigure}{0.45\textwidth}
    \includegraphics[width=\linewidth]{Plots/plt_10_participation_top_10_resnet34_iid.pdf}
  \end{subfigure}
  \caption{Test accuracy vs number of rounds (left) and communicated GB (right) on the Tiny-ImageNet dataset using ResNet-34 with $\gamma{=}0.5$, $p {=} 0.1$.}% across participation rates with Top-10\%
  \label{fig:Tiny-ImageNet_resnet34_100_100_5}
\end{figure*}

\paragraph{IID control experiments.}
To isolate the effect of data heterogeneity, we also evaluate under IID partitions in figure \ref{fig:cifar100_resnet18_iid}. We report \emph{accuracy vs.\ rounds} and \emph{accuracy vs.\ GB} for: (i) partial participation ($q{=}0.5$) with Top-5\% and Top-10\%. Across datasets, SA-PEF matches or exceeds EF/CSER in early rounds under the same communication budget. 

\begin{figure*}[htbp]
  \centering
  % Row 1
  \begin{subfigure}{0.48\textwidth}
    \includegraphics[width=\linewidth]{Plots/efficiency_plt_50_participation_top_5_resnet18_iid.pdf}
    \caption{$p {=} 0.5$, Top-5\%}\label{fig:cifar100_50_5_e_iid}
  \end{subfigure}\hfill
  \begin{subfigure}{0.48\textwidth}
    \includegraphics[width=\linewidth]{Plots/efficiency_plt_50_participation_top_10_resnet18_iid.pdf}
    \caption{$p {=} 0.5$, Top-10\%}\label{fig:cifar100_50_10_e_iid}
  \end{subfigure}
  \vspace{4pt}
  % Row 2
  \begin{subfigure}{0.48\textwidth}
    \includegraphics[width=\linewidth]{Plots/plt_50_participation_top_5_resnet18_iid.pdf}
    \caption{$p {=} 0.5$, Top-5\%}\label{fig:cifar100_50_5_c_iid}
  \end{subfigure}\hfill
  \begin{subfigure}{0.48\textwidth}
    \includegraphics[width=\linewidth]{Plots/plt_10_participation_top_10_resnet9.pdf}
    \caption{$p {=} 0.5$, Top-10\%}\label{fig:cifar100_50_10_c_iid}
  \end{subfigure}
  \caption{Test accuracy vs number of rounds (row 1) and communicated GB (row 2) on the CIFAR-100 dataset using ResNet-18.}% across participation rates with Top-10\%
  \label{fig:cifar100_resnet18_iid}
\end{figure*}

\paragraph{Extreme partial participation and local work.}
To further stress-test our methods, we also consider more demanding FL regimes with very low participation and larger local work. In particular, we run experiments with participation $p=0.05$ and $T=10$ local SGD steps per round, comparing EF, CSER, and SA-PEF under the same compression level. As shown in Fig.~\ref{fig:extreme_pp}, SA-PEF consistently attains higher accuracy than EF and CSER at a fixed communication budget, indicating that its advantages persist even under extreme partial participation and increased local work.

\begin{figure*}[h]
  \centering
  \begin{subfigure}{0.45\textwidth}
    \includegraphics[width=\linewidth]{Plots/efficiency_plt_5_participation_top_1_resnet9.pdf}
  \end{subfigure}\hfill
  \begin{subfigure}{0.45\textwidth}
    \includegraphics[width=\linewidth]{Plots/efficiency_plt_5_participation_top_10_resnet9.pdf}
  \end{subfigure}
   \caption{Test accuracy vs.\ number of rounds on CIFAR-10 with ResNet-9 under Top-1\% (left) and Top-10\% (right) uplink compression.}
  \label{fig:extreme_pp}
\end{figure*}

\paragraph{Wall-clock efficiency.}
To quantify implementation overhead, we report test accuracy versus wall-clock time on CIFAR-10/ResNet-9 under a fixed hardware setup (six NVIDIA RTX A5000 GPUs across 6 nodes) in the Figure~\ref{fig:wall_clock}. SA-PEF reaches a given accuracy level substantially earlier than EF, SAEF, CSER, and FedAvg, confirming that the small extra vector operations it introduces incur negligible runtime cost while improving time-to-accuracy.

\begin{figure*}[h]
  \centering
  \begin{subfigure}{0.45\textwidth}
    \includegraphics[width=\linewidth]{Plots/wall_clock_time_top1.pdf}
  \end{subfigure}\hfill
  \begin{subfigure}{0.45\textwidth}
    \includegraphics[width=\linewidth]{Plots/wall_clock_time.pdf}
  \end{subfigure}
  \caption{Test accuracy vs.\ wall-clock time (s) on CIFAR-10 with ResNet-9 under Top-1\% (left) and Top-10\% (right) uplink compression.}
  \label{fig:wall_clock}
\end{figure*}

\paragraph{Scaled-sign compressor.}
Besides Top-$k$, we also consider a scaled-sign compressor $C(x) = \frac{\|x\|_{1}}{d}\,\mathrm{sign}(x)$. This is the group-scaled sign compressor of \citet{li2023analysis}, specialized to a single block ($M=1$), and Proposition~C.1 in that work shows that it is contractive in the sense of our Definition~1 for a suitable $\delta>0$, so it falls within our theoretical framework. Hence, we report preliminary CIFAR-10/ResNet-9 results with this scaled-sign compressor in Figure~\ref{fig:scaled_sign}, where SA-PEF again achieves higher accuracy than EF, SAEF, and CSER at a fixed communication budget.

\begin{figure*}[h]
  \centering
  \begin{subfigure}{0.45\textwidth}
    \includegraphics[width=\linewidth]{Plots/SignSGD_efficiency_plt_5_participation_top_10_resnet9.pdf}
  \end{subfigure}\hfill
  \begin{subfigure}{0.45\textwidth}
    \includegraphics[width=\linewidth]{Plots/SignSGD_efficiency_plt_5_participation_top_10_resnet9_dir05.pdf}
  \end{subfigure}
  \caption{Test accuracy vs.\ number of rounds on CIFAR-10 with ResNet-9 for Dirichlet-$\gamma$ partitions: $\gamma = 0.1$ (left) and $\gamma = 0.5$ (right).}
  \label{fig:scaled_sign}
\end{figure*}

\paragraph{Effect of momentum.}
To isolate the role of momentum, we repeat our CIFAR-10/100 experiments using SGD \emph{without} momentum, keeping all other hyperparameters and compression settings fixed, and provide the results in Figure~\ref{fig:cifar10and100_resnet9_100_10}. The test-accuracy trajectories and accuracy-communication curves show that SA-PEF consistently matches or outperforms EF, SAEF, and CSER, and remains close to FedAvg in terms of accuracy per communicated
GB. This suggests that our conclusions are essentially momentum-agnostic: momentum slightly reshapes the trajectories but does not drive the gains of SA-PEF over EF-style baselines.

\begin{figure*}[htbp]
  \centering
  % Row 1
  \begin{subfigure}{0.45\textwidth}
    \includegraphics[width=\linewidth]{Plots/momentum_ablation_Cifar10_acc.pdf}
  \end{subfigure}\hfill
  \begin{subfigure}{0.45\textwidth}
    \includegraphics[width=\linewidth]{Plots/momentum_ablation_Cifar100_acc.pdf}
  \end{subfigure}
  \vspace{4pt}
  % Row 2
  \begin{subfigure}{0.45\textwidth}
    \includegraphics[width=\linewidth]{Plots/momentum_ablation_Cifar10_comm.pdf}
    \caption{with CIFAR-10 using ResNet-9}\label{fig:cifar10_10_10_comm}
  \end{subfigure}\hfill
  \begin{subfigure}{0.45\textwidth}
    \includegraphics[width=\linewidth]{Plots/plt_100_participation_top_10_resnet9.pdf}
    \caption{with CIFAR-100 using ResNet-18}\label{fig:cifar100_10_10_cmm}
  \end{subfigure}
  \caption{Test accuracy vs number of rounds (row 1) and communicated GB (row 2) using $p {=} 0.1$, Top-10\%, and $\gamma{=}0.5$.}% across participation rates with Top-10\%
  \label{fig:cifar10and100_resnet9_100_10}
\end{figure*}

\paragraph{Comparison with SCAFCOM.}
To further assess SA-PEF, we follow the MNIST setup of \citet{huang2024stochastic}: a 2-layer fully-connected network is trained on MNIST distributed across $N=200$ clients in a highly heterogeneous regime (each client holds data from at most two classes), with partial participation $p=0.1$ and $10$ local steps. We apply aggressive Top-1\% compression to EF-style methods. As shown in Figure~\ref{fig:mnist_scafcom}, SA-PEF closely tracks SCAFCOM and both substantially outperform standard EF and FedAvg, while uncompressed SCAFFOLD lies in between. This indicates that, under the same communication budget, SA-PEF can match the robustness of SCAFCOM’s control-variate-plus-momentum design while retaining the simpler EF architecture (one residual per client, no additional drift-correction state). 

\begin{figure*}[h]
  \centering
  \begin{subfigure}{0.45\textwidth}
    \includegraphics[width=\linewidth]{Plots/scafcom_sapef_acc.pdf}
    \caption{Test accuracy vs.\ rounds.}
    \label{fig:mnist_scafcom_acc}
  \end{subfigure}\hfill
  \begin{subfigure}{0.45\textwidth}
    \includegraphics[width=\linewidth]{Plots/scafcom_sapef_loss.pdf}
    \caption{Training loss vs.\ rounds.}
    \label{fig:mnist_scafcom_loss}
  \end{subfigure}
  \caption{Comparison with SCAFCOM on MNIST under Top-1\% compression, partial participation $p{=}0.1$, and $10$ local steps.}
  \label{fig:mnist_scafcom}
\end{figure*}

To better align with our main experimental setup, we additionally report a preliminary CIFAR-10/ResNet-9 experiment under aggressive compression in the Figure~\ref{fig:cifar10_scafcom}. We follow the local mini-batch step formulation of \citet{huang2024stochastic} (rather than local epochs) for a fair comparison. We use $K=100$ clients and $R=200$ communication rounds; at each round, the server samples $10$ clients, and each selected client performs $20$ local mini-batch SGD steps on a ResNet-9 model. For all EF-style methods (EF, SAEF, SA-PEF, and SCAFCOM) we apply Top-$1\%$ and Top-$10\%$ sparsification to the uplink updates, while FedAvg communicates dense updates. SCAFCOM uses control-variate and momentum coefficients $\alpha_{\mathrm{sc}} = 0.1$ and $\beta_{\mathrm{sc}} = 0.2$, selected via a small grid search. SA-PEF again behaves competitively with SCAFCOM while clearly improving over EF, SAEF, and FedAvg under the same communication budget. All curves are averaged over five independent runs with different random seeds. The learning-rate combinations that yield the highest test accuracy are listed in Table~\ref{tab:fed_ef_lrs}.

\begin{table}[t]
    \centering
    \setlength{\tabcolsep}{8pt}
    \begin{tabular}{lcc}
        \toprule
        Algorithm & $\eta$ & $\eta_r$  \\
        \midrule
        SA-PEF& $1$      & $10^{-3}$ \\
        SAEF& $1$      & $10^{-3}$ \\
        EF& $1$      & $10^{-3}$ \\
        SCAFCOM& $3$      & $10^{-1}$ \\
        FedAvg& $1$      & $10^{-1}$ \\
        \bottomrule
    \end{tabular}
    \caption{Optimal global ($\eta$) and local ($\eta_r$) learning rate combinations.}
    \label{tab:fed_ef_lrs}
\end{table}

\begin{figure*}[h]
  \centering
  \begin{subfigure}{0.45\textwidth}
    \includegraphics[width=\linewidth]{Plots/scafcom_sapef_acc_cifar10_top1.pdf}
  \end{subfigure}\hfill
  \begin{subfigure}{0.45\textwidth}
    \includegraphics[width=\linewidth]{Plots/scafcom_sapef_acc_cifar10.pdf}
  \end{subfigure}
  \caption{Test accuracy vs.\ number of rounds on CIFAR-10 with ResNet-9, partial participation $p{=}0.1$, and $10$ local steps under Top-1\% (left) and Top-10\% (right) uplink compression.}
  \label{fig:cifar10_scafcom}
\end{figure*}



\end{document}
